\documentclass{article}

\usepackage{iclr2026_conference,times}

\input{math_commands.tex}

\usepackage{amsmath}
\usepackage{amssymb}
\usepackage{amsthm}
\usepackage{mathtools}
\usepackage{hyperref}
\usepackage{url}
\usepackage{flafter}
\usepackage{placeins}
\usepackage[inkscapelatex=false]{svg}


\newcommand{\Prob}{\mathbb{P}}
\newcommand{\diag}{\operatorname{diag}}
\newcommand{\rank}{\operatorname{rank}}
\newcommand{\srank}{\operatorname{srank}}
\newcommand{\erank}{\operatorname{erank}}
\newcommand{\St}{\operatorname{St}}
\newcommand{\op}{\operatorname{op}}
\newcommand{\weak}{\Rightarrow}
\newcommand{\one}{\mathbf{1}}
\newcommand{\hadamard}{\mathbin{\odot}}
\newcommand{\bA}{\mathbf{A}}
\newcommand{\bB}{\mathbf{B}}
\newcommand{\bC}{\mathbf{C}}
\newcommand{\bD}{\mathbf{D}}
\newcommand{\bM}{\mathbf{M}}
\newcommand{\bL}{\mathbf{L}}
\newcommand{\bP}{\mathbf{P}}
\newcommand{\bS}{\mathbf{S}}
\newcommand{\bT}{\mathbf{T}}
\newcommand{\bQ}{\mathbf{Q}}
\newcommand{\bK}{\mathbf{K}}
\newcommand{\bU}{\mathbf{U}}
\newcommand{\bV}{\mathbf{V}}
\newcommand{\be}{\mathbf{e}}
\newcommand{\bu}{\mathbf{u}}
\newcommand{\bv}{\mathbf{v}}
\newcommand{\bw}{\mathbf{w}}
\newcommand{\bG}{\mathbf{G}}
\newcommand{\bH}{\mathbf{H}}
\newcommand{\bSigma}{\boldsymbol{\Sigma}}
\newcommand{\ESD}{\mu_{\mathrm{eig}}}
\newcommand{\SSD}{\mu_{\mathrm{sing}}}
\newcommand{\barPhi}{\overline{\Phi}}

\newtheorem{theorem}{Theorem}
\newtheorem{lemma}[theorem]{Lemma}
\newtheorem{proposition}[theorem]{Proposition}
\newtheorem{corollary}[theorem]{Corollary}
\newtheorem{conjecture}[theorem]{Conjecture}
\newtheorem{assumption}[theorem]{Assumption}
\theoremstyle{definition}
\newtheorem{definition}[theorem]{Definition}
\theoremstyle{remark}
\newtheorem{remark}[theorem]{Remark}

\title{Attention Non-Collapse in State-Space Duality Models}

% Authors must not appear in the submitted version. They should be hidden
% as long as the \iclrfinalcopy macro remains commented out below.
\author{Mengda Li \\
       \And
       Zhaorui Dong \\
       \And
       Jianfeng Yao}

%\iclrfinalcopy

\begin{document}

\maketitle

\begin{abstract}
State-space models are emerging as compelling linear-time alternatives to
Transformers.  Yet efficiency alone does not explain what fundamentally
separates the two architectures: how do they organize and preserve information
across a sequence?  We answer this question through \emph{State Space Duality}
(SSD), which expresses linear attention, Mamba, and RetNet using a common
structured attention matrix,
\(
    \bM=\bL\hadamard\bQ\bK^\top.
\)
Viewed through its spectrum, this matrix reveals a sharp divide.  Softmax
self-attention concentrates information into a few dominant directions, while
the structured causal masks of Mamba and RetNet keep information spread across
many directions as context length grows.  We show that this contrast persists
in models of both learned features and random initialization.  We also identify
the role of memory decay: without it, causal linear attention can form a
dominant direction even though the remaining spectrum does not collapse.  In
particular, self-attention
Transformers are \emph{information compressors}, whereas SSD
models are \emph{information diffusers}.
\end{abstract}

\section{Introduction}
Transformers and state-space models (SSMs) are two prominent approaches to
sequence modeling.  Their computational distinction is well known: dense
self-attention mixes a sequence at quadratic cost, whereas recurrent SSM
mixers can be evaluated in linear time.  Computation alone, however, does not
explain how these architectures distribute signal across token directions.
Recent spectral analyses show that, under isotropic or bounded low-rank score
models, row-softmax normalization can preserve a small number of dominant
directions while most empirical spectral mass vanishes as the sequence grows
\citep{li:2026spectral}.  We refer to this behavior as \emph{spectral
compression}.  Do the structured causal operators used by linear-time models
compress information in the same way, or do they spread it across many
directions?

We investigate this question using the structured masked-attention view
provided by \emph{State Space Duality} (SSD) \citep{mamba2}.  At the level of
token-to-token mixing, several scalar-transition recurrent mechanisms can be
represented as
\(
        \bM=\bL\hadamard \bQ\bK^\top
        \)
where $\bQ\bK^\top$ is a query-key score matrix, $\bL$ is a structured causal
mask, and $\hadamard$ denotes entrywise multiplication.  We focus on the causal
all-one mask of unnormalized linear attention \citep{linear_attn}, the
$1$-semiseparable mask associated with scalar SSD/Mamba transitions
\citep{mamba1,mamba2}, and the exponential Toeplitz mask used by RetNet-style
retention \citep{sun:2023retentive}.  

The Hadamard product in the structured attention can turn a low-rank score matrix into a
full-rank causal operator.  The relevant question is therefore spectral rather
than algebraic: is the resulting mass carried by a few dominant directions or
spread over a growing number of directions?  To answer it precisely, we
distinguish two notions of collapse.  \emph{Absolute spectral degeneration}
means that an unrescaled empirical spectral law converges to $\delta_0$.
\emph{Relative rank collapse} means that a scale-invariant rank proxy remains
bounded as the sequence length grows.  A matrix may exhibit the first behavior
without the second.  We use \emph{spectral diffusion} when normalized spectral
mass remains distributed over many directions, even if the overall scale
shrinks.

Our first regime is a fixed-rank proxy for trained or compressed query and key
matrices.  Because $\bQ=XW_Q$ and $\bK=XW_K$ share the same input, their
token-side singular spaces may be strongly dependent.  We model an extreme
form of this dependence with a common isotropic token frame, while taking the
projection-side singular spaces to be independent and isotropic.  These are
tractable idealizations, not consequences of training.  Under them, every
unit-diagonal lower-triangular mask has a macroscopic effective eigenvalue
support at the natural scale $\sigma_a^Q\sigma_b^K\asymp\ell\sqrt d$.  In the
rank-one case, singular values also reveal the off-diagonal action of the mask:
stable semiseparable and exponential-decay masks have nearly linear rank lower
bounds, whereas the undiscounted causal all-one mask has a much larger top
direction and correspondingly smaller rank proxies.

Our second regime is an independent Gaussian initialization proxy.  The
eigenvalues are then product-normal diagonal inner products for every
unit-diagonal causal mask.  For each fixed projection dimension, their stable
rank is of order $\ell/\log\ell$ and their entropy effective rank is of order
$\ell$.  Singular values again distinguish the masks.  Causal linear attention
has a nondegenerate squared-singular-value limit and near-linear rank growth.
For stable semiseparable and exponential-decay masks, the unrescaled law instead
converges to $\delta_0$, yet all three scale-invariant rank proxies diverge at nearly
linear rates.  This initialization result is the clearest reason to keep
absolute degeneration and relative rank collapse separate throughout the
paper.

\subsection{Related Work}

The starting point is the Transformer self-attention mechanism of
\citet{vaswani:2017attention}.  A growing theoretical literature studies when
attention matrices lose rank or spectral mass.  Rank collapse across depth was
identified by \citet{attention_is_not_depth} and further connected to signal
propagation by \citet{anagnostidis2022signal}.  Collapse in width has also been
studied in an i.i.d.\ Markov-matrix setting \citep{mindthegap_width}, 
while \citet{li:2026spectral} establish degeneration of dense and causal
softmax attention as sequence length grows under isotropic score geometry.
At initialization, \citet{hayase2026gaussian} characterize the asymptotic
spectral distribution of attention through Gaussian equivalence.  Our work
studies the complementary role of structured causal masks and gives direct
low-rank softmax comparisons in Appendix~\ref{app:softmax-comparisons}.

The quadratic sequence-length complexity of self-attention has motivated a
large class of recurrent alternatives, even as implementations such as
FlashAttention substantially reduce its memory traffic \citep{flashattention}.
Linear attention reorganizes causal attention into a recurrent computation
\citep{linear_attn}.  RetNet introduces head-dependent exponential retention
\citep{sun:2023retentive}, while selective SSMs and Mamba use input-dependent
state transitions \citep{mamba1,mamba2}.  State Space Duality places
scalar-transition SSM mixers in the same structured-matrix language as masked
attention \citep{mamba2}.  Toeplitz and Fourier token mixers provide related,
but distinct, structured alternatives \citep{qin:2023toeplitz,leethorp:2022fnet}.

\subsection{Structured Masked-Attention Model}

Let $\ell$ be the sequence length, let $d$ be the projection dimension, and let
$\hadamard$ denote the Hadamard product.  Throughout the paper, causal masks
are lower triangular and unit diagonal.  The token-token interaction matrix is
\begin{equation}
    \bM=\bL\hadamard \bQ\bK^\top,\qquad
    \bQ,\bK\in\R^{\ell\times d},\quad
    \bL\in\R^{\ell\times \ell}.
    \label{eq:model}
\end{equation}
Here $\bQ\bK^\top$ is the query-key score matrix and $\bL$ is a structured
causal mask.  The models studied in this paper correspond to the following
unit-diagonal lower-triangular masks.

\begin{enumerate}
\item \textbf{Causal linear attention} \citep{linear_attn}.  The causal mask is
\[
    \bL_{ij}=\one_{\{i\geq j\}}.
\]
Here $\bM$ is the unnormalized causal score matrix.  This baseline does not
include the feature maps or row normalizers used by some linear-attention
implementations.

\item \textbf{Scalar SSD/Mamba transition} \citep{mamba1,mamba2}.  The
\(1\)-semiseparable mask is
\begin{equation}
    \bL_{ij}
    =
    \one_{\{i\geq j\}}\prod_{t=j+1}^{i}a_t,
    \qquad 0\leq a_t<1,
    \label{eq:semiseparable-mask}
\end{equation}
with the empty product equal to one.  Results that require uniform stability
will additionally assume that there is a deterministic constant $a<1$ such
that $a_t\leq a$ for every $t$.  In the SSM notation, the query and key matrices
\(\bQ,\bK\) correspond to the \(\bC\)- and \(\bB\)-branches.  In a selective
model, the coefficients $a_t$ can depend on the input.  Our low-rank results
may be read conditionally on the realized mask; the initialization results use
the explicit independence proxy in Assumption~\ref{ass:logistic-normal-mask}.

\item \textbf{Retentive Network (RetNet)} \citep{sun:2023retentive}.  For a
single head, the Toeplitz decay mask is
\[
    \bL_{ij}=\one_{\{i\geq j\}}\gamma^{i-j},
    \qquad 0<\gamma<1.
\]
This is a stationary Toeplitz mask and a special case of a stationary
$1$-semiseparable mask.  We suppress the head index on $\gamma$ and absorb any
query/key positional transformation into $\bQ$ and $\bK$.
\end{enumerate}

\subsection{Contributions}
We study the asymptotic spectral behavior of structured attention in SSD models versus self-attention in Transformers. The main contributions are:

\begin{itemize}
    \item \textbf{Attention non-collapse in SSD models.}  To the best of our
    knowledge, we are the first to identify attention rank non-collapse as a
    characteristic phenomenon of structured SSD attention, to formulate it
    spectrally, and to prove it for the structured mixers underlying Mamba/SSD
    and RetNet.  Their effective dimension grows with sequence length rather
    than collapsing to a fixed number of directions.

    \item \textbf{A compressor--diffuser distinction.}  We give a direct
    spectral comparison with Transformer self-attention.  Whereas softmax
    attention concentrates information into a few dominant directions, the
    structured masks of Mamba and RetNet spread it across many.  This suggests
    that SSM-style mixers and self-attention play complementary roles: the
    former diffuse information, while the latter compress it.

    \item \textbf{A theory spanning learned representations and
    initialization.}  We establish this picture under both a low-rank proxy
    for learned query-key representations and an i.i.d.\ Gaussian initialization
    model by characterizing their spectral distributions and effective ranks.
    The analysis also identifies causal linear attention as a boundary case
    and separates spectral shrinkage from genuine rank collapse.
\end{itemize}

\section{Theory of Rank Non-Collapse}


We use the following rank proxies and asymptotic notation to describe the
spectral behavior of the structured attention matrix $\bM$:
\begin{definition}[Rank proxies and asymptotic notation]
For a nonzero matrix $\bA\in\R^{\ell\times\ell}$ with singular values
$s_1(\bA),\ldots,s_\ell(\bA)$, define
\[
    \srank(\bA)
    :=
    \frac{\|\bA\|_F^2}{\|\bA\|_{\op}^2},
    \qquad
    \erank_1(\bA)
    :=
    \exp\left\{-\sum_{j=1}^{\ell}q_j\log q_j\right\},
    \quad
    q_j=\frac{s_j(\bA)}{\sum_k s_k(\bA)},
\]
and
\[
    \erank_2(\bA)
    :=
    \exp\left\{-\sum_{j=1}^{\ell}\widetilde q_j\log\widetilde q_j\right\},
    \qquad
    \widetilde q_j=\frac{s_j(\bA)^2}{\sum_k s_k(\bA)^2},
\]
with $0\log0=0$; for the zero matrix, set these three singular-value proxies
equal to zero.  For eigenvalues
$\lambda_1(\bA),\ldots,\lambda_\ell(\bA)$, define
\[
    \srank_\lambda(\bA)
    :=
    \frac{\sum_i|\lambda_i(\bA)|}{\max_i|\lambda_i(\bA)|},
    \qquad
    \erank_\lambda(\bA)
    :=
    \exp\left\{-\sum_{i=1}^{\ell}p_i\log p_i\right\},
    \quad
    p_i=\frac{|\lambda_i(\bA)|}{\sum_j|\lambda_j(\bA)|}.
\]
The eigenvalue proxies are used when the spectrum is not identically zero.  By
convention, both are set equal to zero when their normalizing spectral mass
vanishes.
We write $X_\ell=O_{\Prob}(a_\ell)$ if $X_\ell/a_\ell$ is tight, and
$X_\ell=\Theta_{\Prob}(a_\ell)$ if $X_\ell/a_\ell$ is bounded above and below
away from zero in probability.  Deterministic constants implicit in
$O(\cdot)$, $\Omega(\cdot)$, and $\Theta(\cdot)$ may depend on fixed mask
parameters but not on $\ell$ or $d$; subscripts such as $O_R(\cdot)$ and
$\Theta_{r,R}(\cdot)$ indicate the allowed parameter dependence.
\label{def:rank-proxies}
\end{definition}

\subsection{Low-Rank Model}

We first study a fixed-rank proxy motivated by trained or compressed query and
key representations.  The model is deliberately stylized: it isolates the
interaction between low-rank scores and a structured causal mask, rather than
asserting that trained query and key matrices are exactly fixed rank.
We write the singular value decompositions as
\begin{equation}
    \begin{aligned}
    \bQ&=\bU_Q\bSigma_Q\bV_Q^\top
        =\sum_{a=1}^{r_Q}\sigma^Q_a\bu^Q_a(\bv^Q_a)^\top,\\
    \bK&=\bU_K\bSigma_K\bV_K^\top
        =\sum_{b=1}^{r_K}\sigma^K_b\bu^K_b(\bv^K_b)^\top.
    \end{aligned}
    \label{eq:svd}
\end{equation}
where $\rank(\bQ)=r_Q$ and $\rank(\bK)=r_K$.  The columns of
$\bU_Q,\bU_K,\bV_Q,\bV_K$ are orthonormal, and
$\bSigma_Q,\bSigma_K$ contain the nonzero singular values.

For linear projections $\bQ=XW_Q$ and $\bK=XW_K$, the shared input $X$
directly couples the token-side geometry.  By contrast, independent random
initialization of $W_Q$ and $W_K$ motivates an isotropic projection-side null
model.  Training need not preserve this independence, so the following
assumption should be understood as a tractable conditional proxy whose fit can
be assessed empirically.
\begin{assumption}[Independent Haar right spaces]
Given the token-side quantities and singular values, the matrices
\[
        \bV_Q\sim \mathrm{Haar}(\St(d,r_Q)),
        \qquad
        \bV_K\sim \mathrm{Haar}(\St(d,r_K))
\]
are independent.
\label{ass:haar}
\end{assumption}

Both \(\bQ\) and \(\bK\) are linear functions of the \emph{same} input matrix
\(X\), so their column spaces are contained in the common ambient space
\(\operatorname{col}(X)\).  This containment alone does not imply that their
left singular vectors coincide.  To isolate the maximally aligned form of
token-side dependence, we impose the following stronger common-frame
idealization.  Throughout this low-rank model, the nonzero singular values are
treated as deterministic parameters; equivalently, the statements may be read
conditionally on them whenever the frame laws below continue to hold.

\begin{assumption}[Common Haar token space]
Let $r:=\max(r_Q,r_K)$.  The left singular vectors used by $\bQ$ and $\bK$
are columns of a common matrix
\[
    \bU=[\bu_1,\ldots,\bu_r]\sim \mathrm{Haar}(\St(\ell,r)),
\]
independent of the right singular spaces.  The ranks, and hence $r$, remain
fixed in all asymptotic statements below.  Thus
\(
    \bu^Q_a=\bu_a,  
    \bu^K_b=\bu_b.
\)
\label{ass:common-left-haar}
\end{assumption}

$\mathrm{Haar}(\St(n,r))$ denotes the uniform Stiefel law, equivalently the
law of the first $r$ columns of an $n\times n$ Haar orthogonal matrix.  The
common Haar frame removes any preferred token-space orientation while retaining
maximal alignment between the left singular vectors of $\bQ$ and $\bK$.  This
is an isotropic modeling choice analogous to random-rotation assumptions in
score-geometry and vector-compression analyses
\citep{RaBitQ24,RaBitQ25,zandieh2026turboquant,li:2026spectral}.

\subsubsection{Eigenvalues of Lower-Triangular Masks}

\begin{theorem}[Eigenvalues of lower-triangular masked attention]
Assume that $\bL$ is lower triangular and unit diagonal and that
Assumptions~\ref{ass:haar} and~\ref{ass:common-left-haar} hold.  Consider any sequence
$(\ell_n,d_n)\to(\infty,\infty)$, suppress the index $n$ from $\ell,d,\bM$, and
suppose
\[
    \theta_{ab}^{(n)}
    :=
    \frac{\sigma^Q_{a,n}\sigma^K_{b,n}}{\ell_n\sqrt {d_n}}
    \longrightarrow \theta_{ab}.
\]
All limits in this theorem and its eigenvalue consequences are along this
joint sequence.
Then, for any fixed collection of distinct tokens $i_1,\ldots,i_m$,
\begin{equation}
    \left(\lambda_{i_1}(\bM),\ldots,\lambda_{i_m}(\bM)\right)
    \weak
    \left(\Lambda_{i_1},\ldots,\Lambda_{i_m}\right),
    \qquad
    \Lambda_i
    :=
    \sum_{a=1}^{r_Q}\sum_{b=1}^{r_K}
    \theta_{ab}\xi_{i,a}\xi_{i,b}G_{ab},
    \label{eq:common-haar-eigenvalue-limit}
\end{equation}
where the variables $\xi_{i,a}$ are independent $N(0,1)$ over token and
singular-vector indices, and $(G_{ab})$ is an independent matrix of independent
$N(0,1)$ entries.  The Gaussian matrix $(G_{ab})$ is shared across all tokens,
whereas the token-side variables $\xi_{i,a}$ depend on $i$.  Define
\(
    D_{\boldsymbol{\theta}}
    :=
    \sum_{a=1}^{\min(r_Q,r_K)}\theta_{aa}^2,
    F_{\boldsymbol{\theta}}
    :=
    \sum_{a=1}^{r_Q}\sum_{b=1}^{r_K}\theta_{ab}^2.
\)
Consequently, the limiting covariance matrix is
\begin{equation}
    \operatorname{Cov}\left(\Lambda_{i_p},\Lambda_{i_q}\right)
    =
    \begin{cases}
    F_{\boldsymbol{\theta}}+2D_{\boldsymbol{\theta}},
        & p=q,\\
    D_{\boldsymbol{\theta}},
        & p\neq q.
    \end{cases}
    \label{eq:common-haar-eigenvalue-covariance}
\end{equation}
\label{thm:lower-triangular-eigenvalues}
\end{theorem}

The natural nontrivial eigenvalue scale in this low-rank model is
$\sigma^Q_a\sigma^K_b\asymp \ell\sqrt d$.  The limiting eigenvalues are
generally dependent because they share the same matrix $\bG$.


\begin{theorem}[Eigenvalue spectral phenomena]
Assume the setting of Theorem~\ref{thm:lower-triangular-eigenvalues},
$\boldsymbol{\theta}\neq0$.  Then the eigenvalues exhibit the
following count and rank phenomena.  Write
$\barPhi(t):=\Prob\{Z>t\}$ for $Z\sim N(0,1)$.

\begin{enumerate}
\item[(i)] Effective eigenvalue counts.  For every $\varepsilon>0$,
\begin{equation}
    \lim_{n\to\infty}
    \frac{\E\#\{i:|\lambda_i(\bM)|>\varepsilon\}}{\ell_n}
    \geq
    \frac{3}{70}
    \barPhi\left(
    \frac{\sqrt2\,\varepsilon}{s_{\boldsymbol{\theta}}}
    \right)>0,
    \label{eq:expected-effective-count}
\end{equation}
where
$s_{\boldsymbol{\theta}}^2=\sum_{a,b}
(1+2\mathbf 1_{\{a=b\}})\theta_{ab}^2$.  The exact limiting count and sharper
tail estimates are stated in Proposition~\ref{prop:eigenvalue-effective-counts}.

\item[(ii)] Eigenvalue stable-rank and effective-rank proxies.  With the
notation of Definition~\ref{def:rank-proxies},
\begin{equation}
    \srank_{\lambda}(\bM)
    =
    \Theta_{\Prob}\left(\frac{\ell}{\log\ell}\right),
    \qquad
    \erank_{\lambda}(\bM)=\Theta_{\Prob}(\ell).
    \label{eq:eigenvalue-rank-summary}
\end{equation}
Sharper random prefactors and the entropy limit are stated in
\eqref{eq:eigenvalue-stable-rank-prefactor} and
\eqref{eq:eigenvalue-entropy-limit}.
\end{enumerate}
\label{thm:eigenvalue-spectral-phenomena}
\end{theorem}

\subsubsection{Singular Values of Unit-Diagonal Causal Masks}

Eigenvalues only see the diagonal of a lower triangular mask.  Singular values
see the full causal operator.  The following rank-one theorem records the
analogues of the eigenvalue phenomena:

\begin{theorem}[Singular-value spectral phenomena]
Assume that $r_Q=r_K=1$ and that Assumptions~\ref{ass:haar}
and~\ref{ass:common-left-haar} hold.  Let $\ell\to\infty$ and
$d=d_\ell\to\infty$
along a joint sequence, and suppose
\[
    \frac{\sigma^Q\sigma^K}{\ell\sqrt d}\longrightarrow \theta\neq0.
\]
Let $\bL$ be one of the following unit-diagonal causal masks: the causal linear
mask $L_{ij}=\one_{\{i\geq j\}}$, a decay mask
$L_{ij}=\one_{\{i\geq j\}}\gamma^{i-j}$ with $0\leq\gamma<1$, or a
$1$-semiseparable mask~\eqref{eq:semiseparable-mask} with
$0\leq a_t\leq a<1$.  Let $C_{\mathrm{inv}}$ be a constant such that every
nonempty principal submatrix $\bL_S$ satisfies
$\|\bL_S^{-1}\|_{\op}\leq C_{\mathrm{inv}}$; for these masks such a constant
depends only on the mask family; see
Proposition~\ref{prop:principal-inverse-bound}.  Let
\(
    \mu_{\mathrm{sing},\ell}
    :=
    \frac{1}{\ell}\sum_{j=1}^{\ell}\delta_{s_j(\bM)}
\)
be the singular-value empirical measure, and use the rank proxies in
Definition~\ref{def:rank-proxies}.  Then the singular values exhibit the
following phenomena.

\begin{enumerate}
\item[(i)] Nondegeneracy of the singular-value empirical law.  Define
\[
    Z_d:=\sqrt d\,\langle\bv^Q,\bv^K\rangle,
    \qquad
    a_{\ell,d}:=\frac{\sigma^Q\sigma^K}{\ell\sqrt d},
\]
so that $Z_d\weak G\sim N(0,1)$ and $a_{\ell,d}\to\theta$.  For every
$t>0$, write $p_t:=\Prob\{|N(0,1)|>t\}=2\barPhi(t)$.  Then, for every
$0<\eta<p_t$,
\begin{equation}
    \Prob\left\{
        \mu_{\mathrm{sing},\ell}
        \left(
            \left(
            \frac{|a_{\ell,d}Z_d|t^2}{2C_{\mathrm{inv}}},
            \infty
            \right)
        \right)
        \geq p_t-\eta
    \right\}
    \longrightarrow1.
    \label{eq:singular-noncollapse}
\end{equation}
Moreover, for every $g_0>0$, set
$\varepsilon_{t,g_0}:=|\theta|g_0t^2/(4C_{\mathrm{inv}})$.  Conditional on
$|Z_d|\geq g_0$,
$\mu_{\mathrm{sing},\ell}((\varepsilon_{t,g_0},\infty))\geq p_t-\eta$
with probability tending to one.  Unconditionally,
\begin{equation}
    \liminf_{\ell\to\infty}
    \Prob\left\{
        \mu_{\mathrm{sing},\ell}
        ((\varepsilon_{t,g_0},\infty))
        \geq p_t-\eta
    \right\}
    \geq \Prob\{|G|\geq g_0\}>0.
    \label{eq:singular-noncollapse-deterministic}
\end{equation}
In particular, $\mu_{\mathrm{sing},\ell}$ does not converge weakly to
$\delta_0$ in probability.

\item[(ii)] Stable rank and effective ranks for decaying masks.  For the decay
mask and the $1$-semiseparable mask,
\begin{equation}
    \srank(\bM)
    \geq
    c\,\frac{\ell}{(\log\ell)^2},
    \qquad
    \erank_{1}(\bM)
    \geq
    c\,\frac{\ell}{\log\ell},
    \qquad
    \erank_{2}(\bM)
    \geq
    c\,\frac{\ell}{(\log\ell)^2}
    \label{eq:stable-mask-singular-ranks}
\end{equation}
with probability tending to one, for a constant $c>0$ depending only on the
mask parameters.

\item[(iii)] Stable rank and effective ranks for causal linear attention.  For
the causal linear mask,
\[
    \srank(\bM)=\Theta_{\Prob}(1).
\]
Moreover, there exists a constant $c>0$ such that
\begin{equation}
    \erank_{1}(\bM)\geq c\log\ell,
    \qquad
    \erank_{2}(\bM)\geq c
    \label{eq:linear-mask-singular-ranks}
\end{equation}
with probability tending to one.
\end{enumerate}
\label{thm:singular-spectral-phenomena}
\end{theorem}

\begin{remark}[Why causal linear attention has smaller rank proxies]
The principle behind the different rates is the size of the causal operator
itself.  For the decay mask and the uniformly stable $1$-semiseparable mask,
the past is discounted, so $\|\bL\|_{\op}=O(1)$; the largest singular direction
is then only logarithmically larger than a typical Haar coordinate effect, and
the rank proxies are almost linear.  For causal linear attention,
$L_{ij}=\one_{\{i\geq j\}}$ accumulates the whole past with no decay, so the
causal operator has norm of order $\ell$.  This large top direction absorbs a
macroscopic amount of singular mass: the theorem still proves
singular-value nondegeneracy, but only logarithmic $\erank_1$ growth and a
constant non-collapse bound for $\erank_2$ are obtained, not quasi-linear
rates.
\label{rem:linear-attention-rank-growth}
\end{remark}

\subsection{Initialization Model}
\label{sec:initialization-model}

The low-rank model above is a useful trained-score proxy.  At initialization,
however, the query and key matrices are not expected to have fixed rank.  The
simplest null model takes their entries to be independent Gaussians at the
usual projection scale.

Following random-matrix models of attention
\citep{mindthegap_width,hayase2026gaussian}, we take $\bQ$ and $\bK$ to be
independent matrices with i.i.d.\ Gaussian entries.  This law is exact in an
idealized linear projection when the input acts orthogonally and $W_Q,W_K$ are
independent Gaussian matrices, by the orthogonal invariance of the Ginibre
ensemble.  Real inputs, normalization layers, and causal convolutions introduce
additional dependence, so Assumption~\ref{ass:gaussian-init} is a null model
rather than a literal description of every SSM implementation.  A scale-mixture
extension is discussed in Appendix~\ref{app:init-details}.

\begin{assumption}[Independent Gaussian query-key initialization]
Fix scale parameters $\sigma_Q,\sigma_K>0$.
The rows $\mathbf q_i^\top$ of $\bQ$ and $\mathbf k_i^\top$ of $\bK$ are
independent over $i$, and
\[
    \mathbf q_i\sim N\left(0,\frac{\sigma_Q^2}{d}\mathbf I_d\right),
    \qquad
    \mathbf k_i\sim N\left(0,\frac{\sigma_K^2}{d}\mathbf I_d\right),
\]
with $\bQ$ independent of $\bK$.  The mask $\bL$ is lower triangular and
unit-diagonal.
\label{ass:gaussian-init}
\end{assumption}

\subsubsection{Eigenvalues at Initialization}


Because $\bM=\bL\hadamard\bQ\bK^\top$ is triangular, its eigenvalues are again
only the diagonal entries.  Thus the mask family does not change the
eigenvalue distribution at initialization as long as $L_{ii}=1$.

\begin{theorem}[Initialization eigenvalue rank proxies]
Under Assumption~\ref{ass:gaussian-init}, the eigenvalues of $\bM$ are
independent diagonal inner products
$\lambda_i(\bM)=\langle\mathbf q_i,\mathbf k_i\rangle$.  In particular, for
each fixed token,
\begin{equation}
    \frac{\sqrt d}{\sigma_Q\sigma_K}\lambda_i(\bM)
    \weak N(0,1)
    \qquad (d\to\infty).
    \label{eq:init-eigenvalue-gaussian-limit}
\end{equation}
For each fixed $d$, as $\ell\to\infty$,
\begin{equation}
    \srank_\lambda(\bM)=\Theta_{\Prob}\left(\frac{\ell}{\log\ell}\right),
    \qquad
    \erank_\lambda(\bM)
    =
    \Theta_{\Prob}(\ell).
    \label{eq:init-fixed-d-ranks}
\end{equation}
\label{thm:init-eigenvalue-proxies}
\end{theorem}

The scalar factor $\sigma_Q\sigma_K/\sqrt d$ cancels from both eigenvalue rank
proxies.
Thus, for each fixed projection dimension, the eigenvalues shrink in absolute
magnitude while remaining distributed across many token directions: the
entropy effective-rank proxy is linear in $\ell$, and the stable-rank proxy
loses only a logarithmic extreme-value factor.  The exact product-normal law,
rank constants, and a Mamba-style scale-mixture proxy are deferred to
Appendix~\ref{app:init-details}.

\subsubsection{Singular Values at Initialization}

Unlike eigenvalues, singular values depend on the off-diagonal entries of the
causal mask.  We first specify stable semiseparable and exponential-decay
proxies for which this dependence remains tractable.

\begin{assumption}[Stable semiseparable and exponential-decay coefficients]
For initialization models with a $1$-semiseparable mask or a Toeplitz-decay
mask, the mask coefficients are independent of $\bQ$ and $\bK$.  In the
$1$-semiseparable case, we use the product formula in
\eqref{eq:semiseparable-mask} and impose exactly one of the following two
alternatives:
\begin{enumerate}
\item[(a)] the coefficients are logistic-normal,
\[
    a_t = \frac{1}{1+\exp(-Z_t)},
    \qquad
    Z_t\stackrel{\mathrm{i.i.d.}}{\sim}N(0,1);
\]
\item[(b)] the coefficients are deterministic and uniformly stable: there
    exists a fixed constant $a_{\max}\in(0,1)$, independent of $\ell$ and $d$,
    such that $0<a_t\leq a_{\max}$ for all $t$.
\end{enumerate}
In the Toeplitz-decay case,
\[
    L_{ij}=\one_{\{i\geq j\}}\gamma^{i-j},
    \qquad
    \gamma\in(0,1),
\]
where $\gamma$ is deterministic and fixed as $\ell,d\to\infty$.
\label{ass:logistic-normal-mask}
\end{assumption}


For the stable semiseparable and RetNet-style initialization proxies considered
here, the unrescaled squared-singular-value limiting spectral distribution
(LSD) is degenerate.

\begin{theorem}[Degenerate initialization LSD for semiseparable and decay masks]
\label{thm:init-degenerate-semiseparable-decay}
Under Assumptions~\ref{ass:gaussian-init}
and~\ref{ass:logistic-normal-mask}, fix positive $\sigma_Q,\sigma_K$ and let $\bL$
be either the $1$-semiseparable mask or the Toeplitz-decay mask.  Assume that
$\ell,d\to\infty$ and $\ell/d\to\tau\in(0,\infty)$.  Let
\[
    \nu_{\ell}^{(2)}
    :=
    \frac1\ell\sum_{j=1}^{\ell}\delta_{s_j(\bM)^2}
\]
be the empirical measure of squared singular values of $\bM$.  Then
\[
    \nu_{\ell}^{(2)}\weak\delta_0
\]
in probability.
\end{theorem}

\begin{theorem}[Initialization rank lower bounds for semiseparable and decay masks]
\label{thm:init-semiseparable-decay-rank-lower-bounds}
Under Assumptions~\ref{ass:gaussian-init}
and~\ref{ass:logistic-normal-mask}, fix positive $\sigma_Q,\sigma_K$ and let $\bL$
be either the $1$-semiseparable mask or the Toeplitz-decay mask.  Assume that
$\ell,d\to\infty$ and $\ell/d\to\tau\in(0,\infty)$.
Then there is a constant $c>0$, depending only on the fixed model and mask
parameters, such that
\[
    \srank(\bM)
    \geq
    c\,\frac{\ell}{(\log\ell)^2},
    \qquad
    \erank_{1}(\bM)
    \geq
    c\,\frac{\ell}{\log\ell},
    \qquad
    \erank_{2}(\bM)
    \geq
    c\,\frac{\ell}{(\log\ell)^2}
\]
with probability tending to one.
\end{theorem}


\begin{remark}[Degenerate LSD versus non-collapsing rank proxies]
Theorem~\ref{thm:init-degenerate-semiseparable-decay} and
Theorem~\ref{thm:init-semiseparable-decay-rank-lower-bounds} concern different
notions of collapse.  The weak convergence
$\nu_{\ell}^{(2)}\weak\delta_0$ is an absolute-scale statement: a typical
unrescaled squared singular value is $o_{\Prob}(1)$.  The stable rank and
effective ranks are instead scale-invariant measures of how the singular mass is
distributed after normalization.  Thus many singular values may all shrink to
zero while remaining spread over many directions.  For example,
$\ell^{-1/2}\mathbf I_\ell$ has empirical squared-singular-value measure
$\delta_{1/\ell}\weak\delta_0$, but stable rank $\ell$.  In the present
initialization model, the diagonal entries already give squared Frobenius mass
of order $\ell/d$ and nuclear mass of order $\ell/\sqrt d$, while the stable
semiseparable and decay masks keep the operator norm at most logarithmically
larger than the single-entry scale.  This is why the unrescaled LSD can degenerate even though
the relative rank proxies remain nearly linear in $\ell$.
\end{remark}


\begin{theorem}[Initialization rank lower bounds for the causal linear mask]
\label{thm:init-causal-linear-rank-lower-bounds}
Under Assumption~\ref{ass:gaussian-init}, fix positive $\sigma_Q,\sigma_K$ and
let $L_{ij}=\one_{\{i\geq j\}}$.  Assume that $\ell,d\to\infty$ and
$\ell/d\to\tau\in(0,\infty)$.  Then there is a constant $c>0$ and, for every
fixed integer $p\geq2$, a constant $c_p>0$ such that
\[
    \srank(\bM)
    \geq
    c_p\ell^{1-1/p},
    \qquad
    \erank_{1}(\bM)
    \geq
    c\ell,
    \qquad
    \erank_{2}(\bM)
    \geq
    c\ell
    \label{eq:init-causal-linear-rank-lower-bounds}
\]
with probability tending to one.  In particular, for every fixed
$\varepsilon>0$,
\[
    \srank(\bM)
    \geq
    c_{\varepsilon}\ell^{1-\varepsilon}
\]
with probability tending to one, and all three rank proxies diverge in
probability.
\end{theorem}


\section{Experiments}
\label{sec:experiments}

\paragraph{Setup.}
We test whether the predicted rank growth appears in trained structured
mixers: layer~0 of Mamba-Codestral-7B and RetNet-1.3B, and layer~1 of
Based-360M, its first global linear-attention block.  We reconstruct the
Mamba-2 SSD matrix, the RetNet decayed causal score matrix, and both the raw and
row-normalized Based Taylor-feature matrices.  At each
$\ell\in\{64,128,256,384,512,768,1024\}$, we independently process two sampled
text spans from the WikiText-103 validation stream.  For each span and each of
four heads, we construct one $\ell\times\ell$ matrix and compute its three rank
proxies; curves show the median of these eight independently computed values.
The same source-text spans are used for every model.  The
semiseparable/decay reference rates for
$(\srank,\erank_1,\erank_2)$ are
$(\ell/(\log\ell)^2,\ell/\log\ell,\ell/(\log\ell)^2)$, while the raw
causal-linear rates are $(1,\log\ell,1)$.  Dashed curves use unit prefactors
because the theorem constants are unknown and are not fitted.  Exact operator
definitions and protocol details are in
Appendix~\ref{app:experimental-details}.

\begin{figure}[t]
  \centering
  \includesvg[width=\linewidth]{figure/structured_rank_comparison_3x3.svg}
  \caption{Rank growth on WikiText-103.  Rows are Based, Mamba-2, and RetNet;
  columns are $\srank$, $\erank_1$, and $\erank_2$.  Lines are medians, dark
  bands are interquartile ranges, and light bands are 10th--90th percentiles.
  Dashed curves are unit-prefactor theoretical rates; the Based rate applies
  only to its raw operator.}
  \label{fig:rank-wikitext}
\end{figure}

\paragraph{Results.}
Mamba-2 has the strongest growth: log--log slopes of median
$(\srank,\erank_1,\erank_2)$ over $\ell\geq256$ are
$(0.839,0.965,0.982)$, so both entropy ranks are nearly linear in the measured
range.  RetNet also diffuses with context, with slopes
$(0.511,0.920,0.708)$; its slower $\srank$ and $\erank_2$ growth indicates a
stronger dominant direction.  In contrast, raw Based has slopes
$(0.028,0.766,0.246)$: stable rank is nearly constant and $\erank_2$ grows
slowly even though $\erank_1$ increases substantially.  Row normalization
changes the levels but not this qualitative separation, giving slopes
$(0.039,0.693,0.239)$.  These trends support the mask-dependent mechanism in
Remark~\ref{rem:linear-attention-rank-growth}: discounted semiseparable/decay
mixers spread normalized singular mass with context, whereas the undiscounted
causal operator retains a persistent top direction.  This is an architectural,
not controlled, comparison across differently trained checkpoints; complete
numerical summaries and limitations are deferred to
Appendix~\ref{app:experimental-details}.

\FloatBarrier

\subsubsection*{Acknowledgments}
Funding and competing-interest disclosures to be added.

\bibliography{ref}
\bibliographystyle{plainnat}

\appendix

\section{Experimental Details}
\label{app:experimental-details}

\paragraph{Operators and layer selection.}
For layer~0 of \texttt{mistralai/Mamba-Codestral-7B-v0.1}, we reconstruct
$\bM_h=\bL_h\hadamard(\bC_h\bB_h^\top)$ with the trained input-dependent
semiseparable mask.  For layer~0 of
\texttt{fla-hub/retnet-1.3B-100B}, we use
$(M_h)_{ij}=\one_{\{i\geq j\}}\gamma_h^{\,i-j}
\langle\mathbf q_{h,i},\mathbf k_{h,j}\rangle/\sqrt{d_h}$, including trained
projections and rotary positions.  For \texttt{hazyresearch/based-360m}, the
raw and normalized operators are
$\bS_h=\operatorname{tril}\{\Phi(\bQ_h)\Phi(\bK_h)^\top\}$ and
$\bP_h=\diag(\bS_h\mathbf 1)^{-1}\bS_h$.

\paragraph{Why Based uses layer~1.}
This is not an indexing offset or a skipped first layer.  Based numbers its
blocks from zero, but the checkpoint uses different mixer types at different
depths.  Its default layer~0 mixer is the local convolutional module
\texttt{BaseConv}, which has neither the Taylor-feature query/key projections
nor the global causal-linear matrix analyzed in Figure~\ref{fig:rank-wikitext}.
The checkpoint configuration sets
\texttt{alt\_mixer\_layers} to $[1,6,11,16,21]$ and defines the corresponding
alternative mixer as Taylor-feature \texttt{LinearAttention}.  Hence layer~1
is the first occurrence of the operator being compared, whereas Mamba-2 and
RetNet already contain their relevant structured mixers at layer~0.  To obtain
the correct trained input to Based layer~1, our extractor applies the token
embedding, the complete layer~0 normalization, \texttt{BaseConv}, residual
connection, MLP, and the layer~1 input normalization before forming the
layer~1 queries and keys.

\paragraph{Sampling and aggregation.}
We concatenate non-heading rows from the raw WikiText-103 validation split and
form one continuous text stream.  At each requested length
$\ell\in\{64,128,256,384,512,768,1024\}$, one source span starts at reference
token~0 and a second starts at a seeded random offset.  Each span is a separate
sample and is forwarded independently through every checkpoint; the
checkpoints receive the same decoded source text, then tokenize it independently
and truncate to $\ell$.  Thus this is a two-sample-per-length evaluation, not
an exhaustive evaluation of every WikiText validation row.  For each sample,
we select four evenly spaced heads and separately construct four
$\ell\times\ell$ matrices, computing all singular values rather than using a
rank approximation.  Each plotted point therefore summarizes $2$ samples
$\times$ $4$ heads $=8$ rank values.  Lines are medians, dark bands are the
25th--75th percentiles, and light bands are the 10th--90th percentiles.

\paragraph{Reference curves.}
The semiseparable and decay rates plotted for
$(\srank,\erank_1,\erank_2)$ are
$(\ell/(\log\ell)^2,\ell/\log\ell,\ell/(\log\ell)^2)$, following
Theorems~\ref{thm:singular-spectral-phenomena}
and~\ref{thm:init-semiseparable-decay-rank-lower-bounds}.  The raw
causal-linear rates are $(1,\log\ell,1)$.  All curves use unit prefactors:
the model-dependent theorem constants are unknown, so their vertical
positions are not numerical tests of the bounds.  The Based reference covers
$\bS_h$ but not the row-normalized $\bP_h$.

\begin{table}[h]
  \centering
  \scriptsize
  {\setlength{\tabcolsep}{3pt}
  \begin{tabular}{lccc}
    \hline
    Operator & $\srank$ & $\erank_1$ & $\erank_2$ \\
    \hline
    Mamba-2
      & $1.99\!\to\!12.04\ (0.839)$
      & $28.23\!\to\!387.04\ (0.965)$
      & $6.94\!\to\!73.31\ (0.982)$ \\
    RetNet
      & $1.43\!\to\!3.51\ (0.511)$
      & $16.98\!\to\!182.57\ (0.920)$
      & $3.43\!\to\!13.74\ (0.708)$ \\
    Based raw
      & $2.06\!\to\!2.51\ (0.028)$
      & $34.78\!\to\!302.42\ (0.766)$
      & $8.84\!\to\!15.17\ (0.246)$ \\
    Based normalized
      & $4.46\!\to\!5.02\ (0.039)$
      & $38.74\!\to\!271.50\ (0.693)$
      & $16.62\!\to\!35.04\ (0.239)$ \\
    \hline
  \end{tabular}}
  \caption{Numerical summary for Figure~\ref{fig:rank-wikitext}.  Each entry
  gives the median at $\ell=64$, the median at $\ell=1024$, and in parentheses
  the least-squares slope of log median rank against $\log\ell$ over
  $\ell\geq256$.}
  \label{tab:trained-rank-summary}
\end{table}

\paragraph{Scope.}
The checkpoints differ in parameter count, training corpus and budget,
projection dimension, and tokenizer, and the experiment samples only four
heads, two windows per length, and one structured layer per model.  Wide
percentile bands, especially for Mamba-2 and RetNet, show meaningful
head-to-head heterogeneity.  Finally, the asymptotic theory uses stylized
low-rank or initialization assumptions that have not been verified for these
trained representations.  Figure~\ref{fig:rank-wikitext} is therefore
architectural evidence for mask-dependent growth patterns, not a controlled
comparison of training recipes or a test of exact asymptotic constants.

\section{Initialization Model Details}
\label{app:init-details}

\subsection{Eigenvalue Distribution}
The proof of Theorem~\ref{thm:init-eigenvalue-proxies} is based on the following
two propositions.  Proposition~\ref{prop:init-eigenvalues} identifies the exact
diagonal product-normal law, and
Proposition~\ref{prop:init-eigenvalue-ranks} converts that law into
stable-rank and entropy-effective-rank estimates.  For clarity, we repeat the
three conclusions used in the main text:
\[
    \frac{\sqrt d}{\sigma_Q\sigma_K}\lambda_i(\bM)
    \weak N(0,1)
    \qquad (d\to\infty),
\]
\[
    \srank_\lambda(\bM)=\Theta_{\Prob}\left(\frac{\ell}{\log\ell}\right),
    \qquad
    \erank_\lambda(\bM)
    =
    \Theta_{\Prob}(\ell)
    \qquad \text{for fixed }d.
\]

\begin{proposition}[Diagonal spectrum at independent Gaussian initialization]
Under Assumption~\ref{ass:gaussian-init}, the eigenvalues of $\bM$ are
independent and
\begin{equation}
    \lambda_i(\bM)
    =
    \langle \mathbf q_i,\mathbf k_i\rangle
    \stackrel{d}{=}
    \frac{\sigma_Q\sigma_K}{d}
    \sum_{a=1}^d Z_{ia}W_{ia},
    \label{eq:init-eigenvalue-product-normal}
\end{equation}
where $Z_{ia},W_{ia}$ are independent standard Gaussians.  Equivalently,
\begin{equation}
    \E e^{t\lambda_i(\bM)}
    =
    \left(
        1-\frac{\sigma_Q^2\sigma_K^2t^2}{d^2}
    \right)^{-d/2},
    \qquad |t|<\frac{d}{\sigma_Q\sigma_K}.
    \label{eq:init-eigenvalue-mgf}
\end{equation}
In particular
\[
    \E\lambda_i(\bM)=0,\qquad
    \operatorname{Var}\{\lambda_i(\bM)\}
    =
    \frac{\sigma_Q^2\sigma_K^2}{d},
\]
and the Gaussian limit \eqref{eq:init-eigenvalue-gaussian-limit} holds.
\label{prop:init-eigenvalues}
\end{proposition}

\begin{proof}
Since $\bL$ is lower triangular and $L_{ii}=1$, the eigenvalues of
$\bM=\bL\hadamard\bQ\bK^\top$ are
\[
    \lambda_i(\bM)=M_{ii}=L_{ii}\langle \mathbf q_i,\mathbf k_i\rangle
    =
    \langle \mathbf q_i,\mathbf k_i\rangle.
\]
The pairs $(\mathbf q_i,\mathbf k_i)$ are independent across $i$, so the
eigenvalues are independent.  Write
$\mathbf q_i=(\sigma_Q/\sqrt d)(Z_{i1},\ldots,Z_{id})$ and
$\mathbf k_i=(\sigma_K/\sqrt d)(W_{i1},\ldots,W_{id})$.  This gives
\eqref{eq:init-eigenvalue-product-normal}.  For independent standard Gaussians
$Z,W$, conditioning on $Z$ gives
\[
    \E\{e^{sZW}\mid Z\}=e^{s^2Z^2/2}.
\]
Averaging over $Z$ yields
\[
    \E e^{sZW}=(1-s^2)^{-1/2},\qquad |s|<1.
\]
Multiplying $d$ independent copies gives \eqref{eq:init-eigenvalue-mgf}.  The
mean and variance follow either by differentiating the moment generating
function at zero or directly from
$\E ZW=0$ and $\operatorname{Var}(ZW)=1$.  Finally,
\[
    \frac{\sqrt d}{\sigma_Q\sigma_K}\lambda_i(\bM)
    =
    \frac{1}{\sqrt d}\sum_{a=1}^d Z_{ia}W_{ia},
\]
so the central limit theorem gives \eqref{eq:init-eigenvalue-gaussian-limit}.
\end{proof}

\begin{proposition}[Eigenvalue rank proxies at initialization]
Let $\srank_\lambda$ and $\erank_\lambda$ be defined as in
Definition~\ref{def:rank-proxies}.  Set
\[
    X_d:=\frac{1}{\sqrt d}\sum_{a=1}^d Z_aW_a,\qquad
    \mu_d:=\E |X_d|,\qquad
    \eta_d:=\E\{|X_d|\log |X_d|\}.
\]
For each fixed $d$, as $\ell\to\infty$,
\begin{equation}
    \srank_\lambda(\bM)=\Theta_{\Prob}\left(\frac{\ell}{\log\ell}\right),
    \qquad
    \frac{\erank_\lambda(\bM)}{\ell}
    \xrightarrow{\Prob}
    c_d:=
    \mu_d\exp\left\{-\frac{\eta_d}{\mu_d}\right\}.
    \label{eq:init-fixed-d-rank-constant}
\end{equation}
\label{prop:init-eigenvalue-ranks}
\end{proposition}

\begin{proof}
By Proposition~\ref{prop:init-eigenvalues},
\[
    \lambda_i(\bM)
    \stackrel{d}{=}
    \frac{\sigma_Q\sigma_K}{\sqrt d}\,X_{i,d},
\]
where $X_{1,d},\ldots,X_{\ell,d}$ are independent copies of $X_d$.  The scalar
factor $\sigma_Q\sigma_K/\sqrt d$ cancels from both $\srank_\lambda$ and
$\erank_\lambda$, so it is enough to analyze the absolute weights
$w_i=|X_{i,d}|$.

For fixed $d$, the moment generating function
$\E e^{tX_d}=(1-t^2/d)^{-d/2}$ is finite for $|t|<\sqrt d$ and singular at
$|t|=\sqrt d$.  Chernoff's bound therefore gives constants $b_d,B_d>0$
such that
\[
    \Prob\{|X_d|>x\}\leq B_de^{-b_dx}
\]
for all sufficiently large $x$.  For the matching lower bound, write
$S_{d-1}=\sum_{a=2}^d Z_aW_a$ (with $S_0=0$), and choose a fixed $D_d<\infty$
such that $\Prob\{|S_{d-1}|\leq D_d\}>0$.  Independence and the standard
product-normal tail bound
$\Prob\{ZW>u\}\geq c u^{-1/2}e^{-u}$ for large $u$ give
\[
\begin{aligned}
\Prob\{X_d>x\}
&\geq
\Prob\{Z_1W_1>\sqrt d\,x+D_d,\ |S_{d-1}|\leq D_d\}\\
&\geq a_d x^{-1/2}e^{-A_dx}
\end{aligned}
\]
for constants $a_d,A_d>0$.  A union bound with the upper tail, and
independence with the lower tail, now show that for constants
$0<c_d<C_d<\infty$,
\[
    \Prob\{c_d\log\ell\leq \max_i w_i\leq C_d\log\ell\}\to1.
\]
The law of large numbers gives
\[
    \frac1\ell\sum_iw_i\to\mu_d
    \qquad\text{and}\qquad
    \frac1\ell\sum_iw_i\log w_i\to\eta_d
\]
in probability; the exponential tail guarantees the needed integrability.
Therefore
\[
    \srank_\lambda(\bM)
    =
    \frac{\sum_iw_i}{\max_iw_i}
    =
    \Theta_{\Prob}\left(\frac{\ell}{\log\ell}\right).
\]
For the entropy effective-rank proxy, use
\[
    -\sum_i p_i\log p_i
    =
    \log\left(\sum_iw_i\right)
    -
    \frac{\sum_iw_i\log w_i}{\sum_iw_i},
    \qquad p_i=\frac{w_i}{\sum_jw_j}.
\]
Substituting the two laws of large numbers gives
\[
    \frac{\erank_\lambda(\bM)}{\ell}
    \xrightarrow{\Prob}
    \mu_d\exp\left\{-\frac{\eta_d}{\mu_d}\right\},
\]
which implies the $\Theta_{\Prob}(\ell)$ statement in
\eqref{eq:init-fixed-d-ranks}.

\end{proof}

\begin{proof}[Proof of Theorem~\ref{thm:init-eigenvalue-proxies}]
Proposition~\ref{prop:init-eigenvalues} gives the diagonal representation and
the Gaussian limit \eqref{eq:init-eigenvalue-gaussian-limit}.  The fixed-$d$
stable-rank estimate and the linear effective-rank estimate follow from
Proposition~\ref{prop:init-eigenvalue-ranks} because
\eqref{eq:init-fixed-d-rank-constant} has a positive finite constant $c_d$.
\end{proof}

\begin{remark}[A Mamba-style initialization proxy]
In SSD/Mamba-2, the query-key factors in the structured-attention view
correspond to the $C$ and $B$ branches.  In common implementations these
branches are produced from the same input by a shared input projection and a
depthwise causal convolution, then split into $B$ and $C$.  A closer
initialization model is therefore not fully i.i.d.\ across tokens.  One useful
intermediate proxy is
\[
    \mathbf q_i\mid \tau_i^Q
    \sim
    N\left(0,\frac{(\tau_i^Q)^2}{d}\mathbf I_d\right),
    \qquad
    \mathbf k_i\mid \tau_i^K
    \sim
    N\left(0,\frac{(\tau_i^K)^2}{d}\mathbf I_d\right),
\]
with $\mathbf q_i$ and $\mathbf k_i$ conditionally independent but with
token-dependent scales inherited from the input, convolution, gate, and
normalization.  To retain the fixed-$d$ rank conclusions above, additionally
assume that, conditional on the entire scale sequence, the row pairs
$(\mathbf q_i,\mathbf k_i)$ are independent across $i$ and are generated by
Gaussian innovations independent of the scales.  Then
\[
    \lambda_i(\bM)\mid \tau_i^Q,\tau_i^K
    \stackrel{d}{=}
    \frac{\tau_i^Q\tau_i^K}{d}
    \sum_{a=1}^d Z_{ia}W_{ia}.
\]
If normalization makes $\tau_i^Q\tau_i^K$ nearly constant, this reduces to
Proposition~\ref{prop:init-eigenvalues}.  If the products
$\tau_i^Q\tau_i^K$ are heterogeneous but uniformly bounded away from zero and
infinity and their empirical law converges in probability to a deterministic
limit, the conditional independence above lets the fixed-$d$ product-normal
calculation yield a stable-rank proxy of order $\ell/\log\ell$ and an
effective-rank proxy linear in $\ell$, with constants adjusted by the scale
mixture.  Varying-$d$ extreme-value constants would require separate
moderate-deviation control.
\label{rem:mamba-init-proxy}
\end{remark}

\subsection{Singular-Value Distribution}

For the causal linear mask, we characterize the nondegenerate
squared-singular-value LSD through its moments.

\begin{theorem}[LSD of causal linear attention]\label{thm:init-singular-causal-mask}
Under Assumption~\ref{ass:gaussian-init}, fix positive $\sigma_Q,\sigma_K$ and
let $L_{ij}=\one_{\{j\leq i\}}$.  Assume that $\ell,d\to\infty$ with
$\ell/d\to\tau\in(0,\infty)$.
Let
\[
    \nu_{\ell}^{(2)}
    :=
    \frac1\ell\sum_{j=1}^{\ell}\delta_{s_j(\bM)^2}
\]
be the empirical measure of squared singular values of $\bM$.  For every fixed
integer $k\ge 1$,
\[
\frac1\ell\operatorname{Tr}\bigl((\bM\bM^\top)^k\bigr)
\longrightarrow
(\sigma_Q\sigma_K)^{2k}m_k(\tau)
\]
in probability, where
\[
    m_k(\tau)
    =
    \sum_{\substack{\pi,\rho\in\mathcal P_2(2k)\\
    |\eta\vee\pi|+|\pi\vee\rho|+|\rho\vee\kappa|=2k+1}}
    \tau^{2k-|\pi\vee\rho|}V_{\pi,\rho}^{\mathrm{caus}}.
\]
Here $\mathcal P_2(2k)$ denotes the set of pair partitions of
$\{1,\ldots,2k\}$,
\[
\kappa=\bigl\{\{2r-1,2r\}:r=1,\ldots,k\bigr\},
\]
and
\[
\eta=\bigl\{\{2r,2r+1\}:r=1,\ldots,k\bigr\},
\qquad 2k+1\equiv 1.
\]
For two partitions $\pi,\rho$, $\pi\vee\rho$ denotes their join.

The quantity $V_{\pi,\rho}^{\mathrm{caus}}\in[0,1]$ is the causal order-polytope
volume defined as follows.  Assign one variable $x_A\in[0,1]$ to each block $A$
of $\eta\vee\pi$, and one variable $y_B\in[0,1]$ to each block $B$ of
$\kappa\vee\rho$.  For each $p=1,2,\ldots,2k$, let $A_p$ denote the block of
$\eta\vee\pi$ containing $p$, and let $B_p$ denote the block of $\kappa\vee\rho$
containing $p$.  Define
\[
V_{\pi,\rho}^{\mathrm{caus}}
=
\int_{[0,1]^{|\eta\vee\pi|+|\kappa\vee\rho|}}
\prod_{r=1}^k
\one_{\{y_{B_{2r-1}}\le x_{A_{2r-1}}\}}
\one_{\{y_{B_{2r}}\le x_{A_{2r}}\}}
\prod_A dx_A\prod_B dy_B.
\]
Consequently, $\nu_{\ell}^{(2)}$ converges weakly in probability to a deterministic
probability measure $\nu_{\tau}^{(2)}$ on $[0,\infty)$, where the fixed scales
$\sigma_Q,\sigma_K$ are suppressed from the notation, satisfying
\[
\int x^k\,d\nu_{\tau}^{(2)}(x)
=
(\sigma_Q\sigma_K)^{2k}m_k(\tau),
\qquad k=1,2,\ldots.
\]
Moreover, $\nu_{\tau}^{(2)}$ has bounded support:
\[
\operatorname{supp}(\nu_{\tau}^{(2)})
\subseteq
\left[0,\sigma_Q^2\sigma_K^2(1+\sqrt\tau)^4\right].
\]
\end{theorem}

\begin{proof}[Proof of Theorem~\ref{thm:init-singular-causal-mask}]
Under Assumption~\ref{ass:gaussian-init}, for the causal mask,
\[
    (\bM)_{ij}
    =
    \one_{\{j\leq i\}}\sum_{a=1}^d q_{ia}k_{ja},
\]
and
\[
    \E q_{ia}q_{i'a'}
    =
    \frac{\sigma_Q^2}{d}\one_{\{i=i'\}}\one_{\{a=a'\}},
    \qquad
    \E k_{ja}k_{j'a'}
    =
    \frac{\sigma_K^2}{d}\one_{\{j=j'\}}\one_{\{a=a'\}}.
\]
With the cyclic convention $i_{k+1}=i_1$,
\[
\begin{aligned}
\frac1\ell\operatorname{Tr}\bigl((\bM\bM^\top)^k\bigr)
&=
\frac1\ell
\sum_{i_1,\ldots,i_k=1}^{\ell}
\sum_{j_1,\ldots,j_k=1}^{\ell}
\sum_{a_1,b_1,\ldots,a_k,b_k=1}^{d}
\prod_{r=1}^k
\one_{\{j_r\leq i_r\}}\one_{\{j_r\leq i_{r+1}\}}
\\
&\qquad\qquad\times
q_{i_ra_r}k_{j_ra_r}q_{i_{r+1}b_r}k_{j_rb_r}.
\end{aligned}
\]
Take expectations and apply Wick's formula separately to the $q_{ia}$-variables
and the $k_{ja}$-variables.  Let the $q$-pairing be $\pi$ and the $k$-pairing be
$\rho$.  The covariance factors contribute
$(\sigma_Q\sigma_K)^{2k}d^{-2k}$.  The feature indices are identified by both
pairings, hence they have $|\pi\vee\rho|$ free classes and contribute
$d^{|\pi\vee\rho|}$.  The $i$-indices are identified by the cyclic trace
identifications $\eta$ and by $\pi$, hence they have $|\eta\vee\pi|$ free
classes.  The $j$-indices are identified by $\kappa$ and by $\rho$, hence they
have $|\rho\vee\kappa|$ free classes.  After these equality constraints, the
causal indicators are exactly the inequalities in the definition of
$V_{\pi,\rho}^{\mathrm{caus}}$, so the Riemann-sum approximation gives
\[
    \ell^{|\eta\vee\pi|+|\rho\vee\kappa|}
    V_{\pi,\rho}^{\mathrm{caus}}
    +
    O\left(\ell^{|\eta\vee\pi|+|\rho\vee\kappa|-1}\right)
\]
admissible choices of the sequence indices.  Therefore, for fixed $k$,
\[
\begin{aligned}
\E\frac1\ell\operatorname{Tr}\bigl((\bM\bM^\top)^k\bigr)
&=
(\sigma_Q\sigma_K)^{2k}
\sum_{\pi,\rho\in\mathcal P_2(2k)}
\left(\frac{\ell}{d}\right)^{2k-|\pi\vee\rho|}
\\
&\qquad\qquad\times
\ell^{|\eta\vee\pi|+|\pi\vee\rho|+|\rho\vee\kappa|-2k-1}
V_{\pi,\rho}^{\mathrm{caus}}
+o(1).
\end{aligned}
\]
We now check which terms can survive.  Put the four pairings
$\eta,\pi,\rho,\kappa$ on the vertex set $\{1,\ldots,2k\}$ and view them as a
ribbon graph with cyclic color order
$\eta\to\pi\to\rho\to\kappa\to\eta$.  The graph is connected because
$|\kappa\vee\eta|=1$.  Its bicolored faces are counted by
$|\eta\vee\pi|$, $|\pi\vee\rho|$, $|\rho\vee\kappa|$, and
$|\kappa\vee\eta|=1$.  Euler's inequality gives
\[
    2k-4k+\bigl(|\eta\vee\pi|+|\pi\vee\rho|+|\rho\vee\kappa|+1\bigr)
    \leq 2,
\]
and hence
\[
    |\eta\vee\pi|+|\pi\vee\rho|+|\rho\vee\kappa|
    \leq
    2k+1.
\]
Thus every exponent of $\ell$ in the expectation is nonpositive, and the
order-one terms are exactly those with equality.  Since $\ell/d\to\tau$,
\[
\E\frac1\ell\operatorname{Tr}\bigl((\bM\bM^\top)^k\bigr)
\longrightarrow
(\sigma_Q\sigma_K)^{2k}
\sum_{\substack{\pi,\rho\in\mathcal P_2(2k)\\
|\eta\vee\pi|+|\pi\vee\rho|+|\rho\vee\kappa|=2k+1}}
\tau^{2k-|\pi\vee\rho|}V_{\pi,\rho}^{\mathrm{caus}}
=
(\sigma_Q\sigma_K)^{2k}m_k(\tau).
\]
It remains to show concentration.  Expanding the second moment gives the same
Wick calculation on two trace cycles, with $4k$ Gaussian positions and
covariance factor $(\sigma_Q\sigma_K)^{4k}d^{-4k}$.  Let
$\widehat\eta,\widehat\kappa\in\mathcal P_2(4k)$ be the disjoint unions of the
two copies of the trace pairings $\eta,\kappa$.  Thus
$|\widehat\kappa\vee\widehat\eta|=2$.  If Wick pairings
$\pi,\rho\in\mathcal P_2(4k)$ connect the two trace cycles, the four-colored
ribbon graph is connected, and Euler's inequality gives
\[
4k-8k+\bigl(|\widehat\eta\vee\pi|+|\pi\vee\rho|
  +|\rho\vee\widehat\kappa|
  +|\widehat\kappa\vee\widehat\eta|\bigr)
\leq 2.
\]
Consequently,
\[
    |\widehat\eta\vee\pi|+|\pi\vee\rho|
    +|\rho\vee\widehat\kappa|
    \leq 4k.
\]
Put $f=|\pi\vee\rho|$ and
$g=|\widehat\eta\vee\pi|+|\rho\vee\widehat\kappa|$.  The causal
inequalities can only reduce the number of sequence-index choices, so a
connected contribution, including the two normalized traces, is
\[
    O\left(d^{f-4k}\ell^{g-2}\right)
    =
    O\left(
      \left(\frac{\ell}{d}\right)^{4k-f}
      \ell^{f+g-4k-2}
    \right)
    =O(\ell^{-2}).
\]
There are only finitely many pairings for fixed $k$.  If the Wick pairings do
not connect the trace cycles, neither pairing has an edge between them, because
$\widehat\eta\vee\widehat\kappa$ already connects all positions within each
trace.  The disconnected Wick sum therefore factors exactly as the square of
the one-trace expectation.  Consequently,
\[
    \E\left(\frac1\ell\operatorname{Tr}\bigl((\bM\bM^\top)^k\bigr)\right)^2
    =
    \left[
    \E\frac1\ell\operatorname{Tr}\bigl((\bM\bM^\top)^k\bigr)
    \right]^2
    +O(\ell^{-2}),
\]
and therefore
\[
    \operatorname{Var}\left(\frac1\ell\operatorname{Tr}\bigl((\bM\bM^\top)^k\bigr)\right)
    =O(\ell^{-2})\longrightarrow 0.
\]
The moment convergence in probability follows from Chebyshev's inequality.

The first moment convergence gives tightness in probability, since for every
positive $R$,
\[
    \nu_{\ell}^{(2)}([R,\infty))
    \leq
    \frac1R\frac1\ell\operatorname{Tr}(\bM\bM^\top).
\]
Set $M_0=1$ and
$M_k=(\sigma_Q\sigma_K)^{2k}m_k(\tau)$ for $k\geq1$.  These numbers form a
Stieltjes moment sequence.  Indeed, for every polynomial
$p(x)=\sum_r c_rx^r$, the expectation limits above imply
\[
    \sum_{r,s}c_rc_sM_{r+s}
    =\lim_{\ell\to\infty}\E\int p(x)^2\,d\nu_\ell^{(2)}(x)\geq0,
\]
and the same argument applied to $xp(x)^2$ gives
$\sum_{r,s}c_rc_sM_{r+s+1}\geq0$.  The Stieltjes moment criterion therefore
gives at least one probability measure on $[0,\infty)$ with these moments.

To see that the moments determine a compactly supported law, compare the
displayed pairing formula with the same Wick formula after deleting the causal
indicators.  Because $0\leq V_{\pi,\rho}^{\mathrm{caus}}\leq1$, the causal
limiting moments are bounded termwise by the corresponding unmasked limiting
moments.  For the unmasked matrix $\bQ\bK^\top$,
\[
    \|\bQ\bK^\top\|_{\op}
    \leq
    \|\bQ\|_{\op}\|\bK\|_{\op},
    \qquad
    \|\bQ\|_{\op}\longrightarrow |\sigma_Q|(1+\sqrt\tau),
    \quad
    \|\bK\|_{\op}\longrightarrow |\sigma_K|(1+\sqrt\tau),
\]
in probability, with the usual fixed-moment bounds for Gaussian operator norms.
Hence
\[
    0\leq
    (\sigma_Q\sigma_K)^{2k}m_k(\tau)
    \leq
    \left(\sigma_Q^2\sigma_K^2(1+\sqrt\tau)^4\right)^k.
\]
Let $\nu_{\tau}^{(2)}$ be any representing probability measure.  For any
positive $\varepsilon$,
\[
    \nu_{\tau}^{(2)}\left(
    \left[\sigma_Q^2\sigma_K^2(1+\sqrt\tau)^4+\varepsilon,\infty\right)
    \right)
    \leq
    \left(
    \frac{\sigma_Q^2\sigma_K^2(1+\sqrt\tau)^4}
    {\sigma_Q^2\sigma_K^2(1+\sqrt\tau)^4+\varepsilon}
    \right)^k.
\]
Letting $k\to\infty$ gives zero mass above
$\sigma_Q^2\sigma_K^2(1+\sqrt\tau)^4+\varepsilon$, and then letting
$\varepsilon$ be arbitrary gives
\[
    \operatorname{supp}(\nu_{\tau}^{(2)})
    \subseteq
    \left[0,\sigma_Q^2\sigma_K^2(1+\sqrt\tau)^4\right].
\]
The same argument applies to every representing measure.  Hence every such
measure is compactly supported in the displayed interval, which gives moment
determinacy and makes $\nu_{\tau}^{(2)}$ unique.

It remains to pass the moments through weak subsequential limits.  Starting
from any subsequence, use the convergence in probability of all integer
moments and a diagonal extraction to choose a further subsequence along which
all of them converge almost surely.  On this probability-one event, convergence
of the first moments gives tightness.  Moreover, for every fixed $k\geq1$,
convergence of the $(k+1)$-st moments gives a uniform bound along this
subsequence and
\[
    \int_{[R,\infty)}x^k\,d\nu_{\ell}^{(2)}(x)
    \leq
    \frac1R\int x^{k+1}\,d\nu_{\ell}^{(2)}(x).
\]
Thus the $k$-th moments are uniformly integrable, so every weak subsequential
limit has moments $(M_k)_{k\geq0}$ and must equal $\nu_{\tau}^{(2)}$.  Every
subsequence therefore has a further subsequence converging almost surely to
$\nu_{\tau}^{(2)}$, which is equivalent to weak convergence in probability of
$\nu_{\ell}^{(2)}$ to $\nu_{\tau}^{(2)}$.
\end{proof}

\begin{proof}[Proof of Theorem~\ref{thm:init-degenerate-semiseparable-decay}]
It suffices to show that
\[
    \frac1\ell\operatorname{Tr}(\bM\bM^\top)
    =
    \frac1\ell\sum_{j=1}^{\ell}s_j(\bM)^2
    \xrightarrow{\Prob}0.
\]
Indeed, for every fixed $\varepsilon>0$,
\[
    \nu_{\ell}^{(2)}([\varepsilon,\infty))
    \leq
    \frac1{\varepsilon\ell}\operatorname{Tr}(\bM\bM^\top).
\]
Conditioning on the mask and using Assumption~\ref{ass:gaussian-init},
\[
    \E\left[\frac1\ell\operatorname{Tr}(\bM\bM^\top)\mid \bL\right]
    =
    \frac{\sigma_Q^2\sigma_K^2}{d\ell}\|\bL\|_F^2.
\]
Thus it remains to check that $\ell^{-1}\E\|\bL\|_F^2$ is bounded uniformly in
$\ell$ for the two masks.

For the $1$-semiseparable mask, we handle the two alternatives in
Assumption~\ref{ass:logistic-normal-mask}.  In the logistic-normal case,
independence of the $a_t$'s gives
\[
\begin{aligned}
\E\|\bL\|_F^2
&=
\sum_{1\leq j\leq i\leq \ell}
\E\prod_{t=j+1}^{i}a_t^2
=
\sum_{r=0}^{\ell-1}(\ell-r)(\E a_1^2)^r
\leq
\frac{\ell}{1-\E a_1^2}.
\end{aligned}
\]
Here $\E a_1^2<1$ because $a_1=(1+\exp(-Z_1))^{-1}\in(0,1)$ almost
surely.  In the deterministic uniformly stable case,
\[
\begin{aligned}
\|\bL\|_F^2
&=
\sum_{1\leq j\leq i\leq \ell}
\prod_{t=j+1}^{i}a_t^2
\leq
\sum_{r=0}^{\ell-1}(\ell-r)a_{\max}^{2r}
\leq
\frac{\ell}{1-a_{\max}^2}.
\end{aligned}
\]
For the Toeplitz-decay mask, the fixed constant $\gamma\in(0,1)$ gives the
deterministic bound
\[
    \|\bL\|_F^2
    =
    \sum_{r=0}^{\ell-1}(\ell-r)\gamma^{2r}
    \leq
    \frac{\ell}{1-\gamma^2}.
\]
Therefore, in either mask family,
\[
    \E\frac1\ell\operatorname{Tr}(\bM\bM^\top)
    =
    \frac{\sigma_Q^2\sigma_K^2}{d\ell}\E\|\bL\|_F^2
    \longrightarrow0.
\]
Markov's inequality gives
$\ell^{-1}\operatorname{Tr}(\bM\bM^\top)\xrightarrow{\Prob}0$, and the first
display implies $\nu_{\ell}^{(2)}\weak\delta_0$ in probability.  Finally, for
the empirical measure of singular values, use
\[
    \frac1\ell\#\{j:s_j(\bM)\geq\varepsilon\}
    \leq
    \frac1{\varepsilon^2\ell}\operatorname{Tr}(\bM\bM^\top)
    \xrightarrow{\Prob}0.
\]
Thus the singular-value LSD also converges to $\delta_0$.
\end{proof}

\begin{proof}[Proof of Theorem~\ref{thm:init-semiseparable-decay-rank-lower-bounds}]
The proof follows the same min-entropy mechanism as parts (ii)--(iii) of
Theorem~\ref{thm:singular-spectral-phenomena}.  We first show that, for both
mask families,
\begin{equation}
    \|\bM\|_{\op}
    \leq
    C\,\frac{\sigma_Q\sigma_K\log\ell}{\sqrt d}
    \label{eq:init-stable-mask-op-bound}
\end{equation}
with probability tending to one.

We start with a row-sum and column-sum bound for the mask:
\begin{equation}
    \max_i\sum_j |L_{ij}|
    \vee
    \max_j\sum_i |L_{ij}|
    \leq
    C\log\ell.
    \label{eq:init-stable-mask-row-column-bound}
\end{equation}
For the Toeplitz-decay mask, the fixed constant $\gamma\in(0,1)$ gives
\[
    \max_i\sum_j |L_{ij}|
    \vee
    \max_j\sum_i |L_{ij}|
    \leq
    \sum_{r=0}^{\ell-1}\gamma^r
    \leq
    \frac1{1-\gamma}
    \leq
    C\log\ell
\]
for all sufficiently large $\ell$, where $C$ may depend on the fixed value of
$\gamma$.  For the $1$-semiseparable mask with deterministic uniformly stable
coefficients,
\[
    \max_i\sum_j |L_{ij}|
    \vee
    \max_j\sum_i |L_{ij}|
    \leq
    \sum_{r=0}^{\ell-1}a_{\max}^r
    \leq
    \frac1{1-a_{\max}}
    \leq
    C\log\ell.
\]
For the $1$-semiseparable mask with logistic-normal coefficients, let $R_\ell$
be the longest consecutive run of indices $t$ for which $a_t>1/2$.  Since the
indicators $\one_{\{a_t>1/2\}}$ are i.i.d. Bernoulli$(1/2)$, a union bound gives
$R_\ell=O(\log\ell)$ with probability tending to one.  On this event, every
block of $R_\ell+1$ consecutive coefficients contains at least one coefficient
at most $1/2$.  Hence, for every $i$ and every $r\geq0$,
\[
    \prod_{t=i-r+1}^{i}a_t
    \leq
    2^{-\lfloor r/(R_\ell+1)\rfloor},
\]
with the obvious convention when the product is empty or the index range is
outside $\{1,\ldots,\ell\}$.  Therefore every row sum and every column sum is at
most
\[
    \sum_{r=0}^{\ell-1}2^{-\lfloor r/(R_\ell+1)\rfloor}
    \leq
    2(R_\ell+1)
    \leq
    C\log\ell,
\]
which proves \eqref{eq:init-stable-mask-row-column-bound} in all cases.

The standard chi-square bound and $\ell/d\to\tau$ give
$\max_i\|\mathbf q_i\|\leq C\sigma_Q$ and
$\max_j\|\mathbf k_j\|\leq C\sigma_K$ with probability tending to one.  On the
event \eqref{eq:init-stable-mask-row-column-bound}, fix $i$ and condition only on
$\bL$ and $\mathbf q_i$; in particular, do not condition on the norms of the
$\mathbf k_j$'s.  The variables
$\langle \mathbf q_i,\mathbf k_j\rangle$ are then independent centered Gaussians
with variance
\[
    \frac{\sigma_K^2\|\mathbf q_i\|^2}{d}
    \leq
    C\frac{\sigma_Q^2\sigma_K^2}{d}
\]
whenever $\|\mathbf q_i\|\leq C\sigma_Q$.  Therefore, for any weights
$0\leq b_j\leq1$ measurable with respect to $\bL$ and satisfying
$\sum_j b_j\leq C\log\ell$, Bernstein's inequality for weighted half-normal
variables gives
\[
    \Prob\left\{
    \sum_j b_j |\langle \mathbf q_i,\mathbf k_j\rangle|
    >
    C\frac{\sigma_Q\sigma_K\log\ell}{\sqrt d}
    \,\middle|\, \bL,\mathbf q_i
    \right\}
    \leq
    \ell^{-3}.
\]
A union bound over $i$, followed by removal of the conditioning and intersection
with the maximum-$\mathbf q_i$ norm event, gives the corresponding bound for all
row sums of $|\bM|$.  For columns, condition only on $\bL$ and $\mathbf k_j$ and
intersect with the maximum-$\mathbf k_j$ norm event; the same argument and a
union bound over $j$ give the identical column-sum bound.  Schur's test then
proves \eqref{eq:init-stable-mask-op-bound}.

Next, because both masks are unit diagonal,
$M_{ii}=\langle \mathbf q_i,\mathbf k_i\rangle$.  The diagonal variables are
independent and satisfy
\[
    \frac{\sqrt d}{\sigma_Q\sigma_K}M_{ii}\weak N(0,1),
    \qquad
    \E\left(\frac{\sqrt d}{\sigma_Q\sigma_K}M_{ii}\right)^2=1,
\]
with uniformly bounded fourth moments.  Thus the triangular-array law of large
numbers gives
\begin{equation}
    \sum_{i=1}^{\ell}|M_{ii}|
    \geq
    c\frac{\sigma_Q\sigma_K\ell}{\sqrt d},
    \qquad
    \sum_{i=1}^{\ell}M_{ii}^2
    \geq
    c\frac{\sigma_Q^2\sigma_K^2\ell}{d}
    \label{eq:init-diagonal-mass-lower-bound}
\end{equation}
with probability tending to one.  The second inequality implies
\[
    \|\bM\|_F^2
    \geq
    \sum_{i=1}^{\ell}M_{ii}^2
    \geq
    c\frac{\sigma_Q^2\sigma_K^2\ell}{d}.
\]
Combining this with \eqref{eq:init-stable-mask-op-bound} gives
\[
    \srank(\bM)
    =
    \frac{\|\bM\|_F^2}{\|\bM\|_{\op}^2}
    \geq
    c\frac{\ell}{(\log\ell)^2}.
\]

For $\erank_1$, use the dual characterization of the nuclear norm.  If the
signs on the diagonal are chosen to match the signs of $M_{ii}$, the resulting
diagonal sign matrix has operator norm at most one.  Hence
\[
    \|\bM\|_*
    \geq
    \sum_{i=1}^{\ell}|M_{ii}|
    \geq
    c\frac{\sigma_Q\sigma_K\ell}{\sqrt d}.
\]
Since Shannon entropy dominates min-entropy,
\[
    \erank_1(\bM)
    \geq
    \frac{\|\bM\|_*}{\|\bM\|_{\op}}
    \geq
    c\frac{\ell}{\log\ell}.
\]
Applying the same min-entropy bound to the squared singular-value weights gives
\[
    \erank_2(\bM)
    \geq
    \frac{\|\bM\|_F^2}{\|\bM\|_{\op}^2}
    =
    \srank(\bM)
    \geq
    c\frac{\ell}{(\log\ell)^2}.
\]
This proves the three claimed lower bounds.
\end{proof}

\begin{proof}[Proof of Theorem~\ref{thm:init-causal-linear-rank-lower-bounds}]
Write $s_j=s_j(\bM)$, and fix an integer $p\geq2$.  By
Theorem~\ref{thm:init-singular-causal-mask}, for every fixed integer $r\geq1$,
\[
    \frac1\ell\sum_{j=1}^{\ell}s_j^{2r}
    =
    \frac1\ell\operatorname{Tr}\bigl((\bM\bM^\top)^r\bigr)
    \xrightarrow{\Prob}
    (\sigma_Q\sigma_K)^{2r}m_r(\tau).
\]
For $r=1$, the limit is positive; indeed $m_1(\tau)=\tau/2$, so
\[
    (\sigma_Q\sigma_K)^2m_1(\tau)
    =
    \frac{\sigma_Q^2\sigma_K^2\tau}{2}
    >
    0.
\]
The limits for $r=1,2,p$ are finite.  Hence we may choose deterministic
constants $0<a<B<\infty$, $C<\infty$, and $C_p<\infty$ such that
\[
    0<a<(\sigma_Q\sigma_K)^2m_1(\tau)<B,
    \qquad
    C>(\sigma_Q\sigma_K)^4m_2(\tau),
    \qquad
    C_p>(\sigma_Q\sigma_K)^{2p}m_p(\tau).
\]
Convergence in probability gives
\[
    \Prob\left\{
    a\ell\leq \sum_{j=1}^{\ell}s_j^2\leq B\ell
    \right\}\to1,
\]
and
\[
    \Prob\left\{
    \sum_{j=1}^{\ell}s_j^4\leq C\ell
    \right\}\to1,
    \qquad
    \Prob\left\{
    \sum_{j=1}^{\ell}s_j^{2p}\leq C_p\ell
    \right\}\to1.
\]
Since this is a finite intersection of high-probability events, with probability
tending to one,
\begin{equation}
    a\ell\leq \sum_{j=1}^{\ell}s_j^2\leq B\ell,
    \qquad
    \sum_{j=1}^{\ell}s_j^4\leq C\ell,
    \qquad
    \sum_{j=1}^{\ell}s_j^{2p}\leq C_p\ell.
    \label{eq:init-causal-linear-moment-event}
\end{equation}

On the event \eqref{eq:init-causal-linear-moment-event},
\[
    \|\bM\|_{\op}^2
    =
    \max_j s_j^2
    \leq
    \left(\sum_{j=1}^{\ell}s_j^{2p}\right)^{1/p}
    \leq
    C_p^{1/p}\ell^{1/p}.
\]
Therefore
\[
    \srank(\bM)
    =
    \frac{\|\bM\|_F^2}{\|\bM\|_{\op}^2}
    =
    \frac{\sum_{j=1}^{\ell}s_j^2}{\|\bM\|_{\op}^2}
    \geq
    c_p\ell^{1-1/p}.
\]

Next we prove the $\erank_2$ bound.  For any probability vector
$(p_j)_j$,
\[
    -\sum_j p_j\log p_j
    \geq
    -\log\left(\sum_j p_j^2\right).
\]
Applying this to
\[
    p_j=\frac{s_j^2}{\sum_{k=1}^{\ell}s_k^2}
\]
gives
\[
    \erank_2(\bM)
    \geq
    \frac{\left(\sum_{j=1}^{\ell}s_j^2\right)^2}
    {\sum_{j=1}^{\ell}s_j^4}
    \geq
    \frac{a^2\ell^2}{C\ell}
    \geq
    c\ell.
\]

It remains to prove the $\erank_1$ bound.  Interpolating the empirical
$L^1$, $L^2$, and $L^4$ norms of the singular values gives
\[
    \frac1\ell\sum_{j=1}^{\ell}s_j
    \geq
    \frac{
    \left(\ell^{-1}\sum_{j=1}^{\ell}s_j^2\right)^{3/2}
    }{
    \left(\ell^{-1}\sum_{j=1}^{\ell}s_j^4\right)^{1/2}
    }
    \geq
    c.
\]
Thus
\[
    \sum_{j=1}^{\ell}s_j\geq c\ell.
\]
Using the same entropy lower bound, now with
\[
    q_j=\frac{s_j}{\sum_{k=1}^{\ell}s_k},
\]
we obtain
\[
    \erank_1(\bM)
    \geq
    \frac{\left(\sum_{j=1}^{\ell}s_j\right)^2}
    {\sum_{j=1}^{\ell}s_j^2}
    \geq
    \frac{c\ell^2}{B\ell}
    \geq
    c\ell.
\]
This proves
\[
    \srank(\bM)
    \geq
    c_p\ell^{1-1/p},
    \qquad
    \erank_1(\bM)\geq c\ell,
    \qquad
    \erank_2(\bM)\geq c\ell
\]
with probability tending to one.  Finally, for any fixed $\varepsilon>0$,
choose an integer $p\geq\max\{2,1/\varepsilon\}$.  Then
$\ell^{1-1/p}\geq \ell^{1-\varepsilon}$, after adjusting the constant, so
\[
    \srank(\bM)\geq c_{\varepsilon}\ell^{1-\varepsilon}
\]
with probability tending to one.
\end{proof}

\section{Low-Rank Model Details}

Write
\begin{equation}
    \bH:=\bV_Q^\top\bV_K\in\R^{r_Q\times r_K}
    \label{eq:H}
\end{equation}
for the right-space overlap matrix.

\begin{lemma}[Haar overlap scale]
Under Assumption~\ref{ass:haar}, for deterministic
$\mathbf a\in\R^{r_Q}$ and $\mathbf b\in\R^{r_K}$,
\[
    \mathbf a^\top\bH\mathbf b
    \stackrel{d}{=}
    \|\mathbf a\|\,\|\mathbf b\|\,\xi_d,
\]
where $\xi_d$ is the inner product of two independent uniform unit vectors in
$\R^d$.  In particular,
\[
    \E\,\mathbf a^\top\bH\mathbf b=0,
    \qquad
    \E|\mathbf a^\top\bH\mathbf b|^2
    =
    \frac{\|\mathbf a\|^2\|\mathbf b\|^2}{d}.
\]
Moreover, for a universal constant $c>0$,
\[
    \Prob\left(
    |\mathbf a^\top\bH\mathbf b|
    >
    t\|\mathbf a\|\|\mathbf b\|
    \right)
    \leq 2\exp(-cdt^2),
    \qquad 0\leq t\leq 1.
\]
If $r_Q,r_K$ are fixed and $d\to\infty$, then
\[
        \sqrt d\,\bH\weak \bG,
\]
where $\bG\in\R^{r_Q\times r_K}$ has independent $N(0,1)$ entries.
\label{lem:haar-overlap}
\end{lemma}

\begin{proof}
By right invariance of the Haar measure on the Stiefel manifold, for every
deterministic vector $\mathbf a\in\R^{r_Q}$, the vector $\bV_Q\mathbf a$ is
uniform on the sphere of radius $\|\mathbf a\|$ in the column space sampled by
$\bV_Q$.  Since the column space is itself uniformly distributed in $\R^d$,
$\bV_Q\mathbf a$ is a uniform vector on the sphere
$\|\mathbf a\|S^{d-1}$.  Similarly, $\bV_K\mathbf b$ is uniform on
$\|\mathbf b\|S^{d-1}$, and the two vectors are independent.  Therefore
\[
    \mathbf a^\top\bH\mathbf b
    =
    \langle \bV_Q\mathbf a,\bV_K\mathbf b\rangle
\]
has the same distribution as
$\|\mathbf a\|\,\|\mathbf b\|\,\xi_d$, where $\xi_d$ is the first coordinate of
a uniform vector on $S^{d-1}$.  Thus $\E\xi_d=0$ and $\E\xi_d^2=1/d$, giving the
displayed second moment.  The concentration bound is the standard spherical-cap
estimate for a coordinate of a uniform vector on the sphere.

It remains to justify the joint Gaussian limit.  Let
$\mathbf Y_Q\in\R^{d\times r_Q}$ and $\mathbf Y_K\in\R^{d\times r_K}$ have independent
standard Gaussian entries.  The QR representation of Haar Stiefel matrices gives
\[
    \bV_Q
    \stackrel{d}{=}
    \mathbf Y_Q(\mathbf Y_Q^\top\mathbf Y_Q)^{-1/2},
    \qquad
    \bV_K
    \stackrel{d}{=}
    \mathbf Y_K(\mathbf Y_K^\top\mathbf Y_K)^{-1/2}.
\]
Consequently,
\[
    \sqrt d\,\bH
    \stackrel{d}{=}
    \left(\frac{\mathbf Y_Q^\top\mathbf Y_Q}{d}\right)^{-1/2}
    \left(\frac{\mathbf Y_Q^\top\mathbf Y_K}{\sqrt d}\right)
    \left(\frac{\mathbf Y_K^\top\mathbf Y_K}{d}\right)^{-1/2}.
\]
Since $r_Q,r_K$ are fixed,
$d^{-1}\mathbf Y_Q^\top\mathbf Y_Q\to\mathbf I_{r_Q}$ and
$d^{-1}\mathbf Y_K^\top\mathbf Y_K\to\mathbf I_{r_K}$ in probability.  Conditional on
$\mathbf Y_Q$, the middle matrix is Gaussian with independent columns and row
covariance $d^{-1}\mathbf Y_Q^\top\mathbf Y_Q$, hence it converges jointly to a matrix with
independent $N(0,1)$ entries.  Slutsky's theorem proves
$\sqrt d\,\bH\weak\bG$.
\end{proof}

\begin{lemma}[Token-side Haar products]
Under Assumption~\ref{ass:common-left-haar}, for any fixed token $i$ and fixed
indices $a,b$,
\[
    \E\,u_a(i)u_b(i)=\delta_{ab}\,\mathrm{Wg}^{O(\ell)}(1)
    =
    \frac{\delta_{ab}}{\ell}.
\]
Moreover,
\[
    \E\,u_a(i)^2u_b(i)^2
    =
    \begin{cases}
    \displaystyle \frac{3}{\ell(\ell+2)}, & a=b,\\[1ex]
    \displaystyle \frac{1}{\ell(\ell+2)}, & a\neq b.
    \end{cases}
\]
Equivalently, for fixed $m,r$ and distinct fixed tokens $i=1,\ldots,m$,
\[
    \left(\sqrt{\ell}\,(u_1(i),\ldots,u_r(i))\right)_{i=1}^m
    \weak
    \left((\xi_{i,1},\ldots,\xi_{i,r})\right)_{i=1}^m,
    \qquad
    \xi_{i,a}\stackrel{\mathrm{i.i.d.}}{\sim}N(0,1),
\]
and hence
\[
    \left(\ell\,u_a(i)u_b(i)\right)_{i,a,b}
    \weak
    \left(\xi_{i,a}\xi_{i,b}\right)_{i,a,b}.
\]
\label{lem:token-haar-products}
\end{lemma}

\begin{proof}
Let $O\in O(\ell)$ be Haar.  The first $r$ columns of $O$ have the same law as
$\bU$.  We apply the orthogonal Weingarten formula of
\citet{collins:2006integration}.  For two entries, the only pairing pairs the
two column indices, so
\[
    \E O_{ia}O_{ib}
    =
    \delta_{ab}\,\mathrm{Wg}^{O(\ell)}(1)
    =
    \frac{\delta_{ab}}{\ell}.
\]
For four entries, write the column-index word as $(a,a,b,b)$.  The orthogonal
Weingarten formula gives
\[
    \E O_{ia}O_{ia}O_{ib}O_{ib}
    =
    \sum_{\pi,\sigma\in\mathcal P_2(4)}
    \delta_i^\pi\,\delta_{aabb}^{\sigma}\,
    \mathrm{Wg}^{O(\ell)}(\pi^{-1}\sigma),
\]
where $\mathcal P_2(4)$ is the set of pair partitions of four points.  Since all
row indices are equal, $\delta_i^\pi=1$ for every $\pi$.  If $a\neq b$, only
the pairing $\{\{1,2\},\{3,4\}\}$ is compatible with the column word; if $a=b$,
all three pairings are compatible.  Equivalently, the squared row
$(O_{i1}^2,\ldots,O_{i\ell}^2)$ has the Dirichlet distribution with all
parameters $1/2$.  Therefore
\[
    \E O_{ia}^2O_{ib}^2
    =
    \frac{1+2\delta_{ab}}{\ell(\ell+2)}.
\]
This is exactly the displayed formula.

For the finite-row Gaussian limit, use the Gaussian representation
$\bU=\mathbf Y(\mathbf Y^\top\mathbf Y)^{-1/2}$ with
$\mathbf Y\in\R^{\ell\times r}$ standard Gaussian.  If $\mathbf y_i$ is the $i$th row,
then
\[
    \sqrt\ell\,(u_1(i),\ldots,u_r(i))
    =
    \mathbf y_i
    \left(\frac{\mathbf Y^\top\mathbf Y}{\ell}\right)^{-1/2}.
\]
For fixed $m,r$, the rows $\mathbf y_1,\ldots,\mathbf y_m$ are independent
$N(0,\mathbf I_r)$ and $\ell^{-1}\mathbf Y^\top\mathbf Y\to\mathbf I_r$ in probability.
Slutsky's theorem gives the joint Gaussian limit for the first $m$ token rows.
The convergence of the products
$\ell u_a(i)u_b(i)$ follows by the continuous mapping theorem.
\end{proof}

\begin{lemma}[Empirical Haar-row law of large numbers]
Let $r$ be fixed and let
$\bU_\ell=[\bu_1,\ldots,\bu_r]\sim\mathrm{Haar}(\St(\ell,r))$.  Define
\[
    \mathbf z_i
    :=
    \sqrt\ell\,(u_1(i),\ldots,u_r(i))\in\R^r,
    \qquad
    \mu_\ell
    :=
    \frac1\ell\sum_{i=1}^{\ell}\delta_{\mathbf z_i}.
\]
Then $\mu_\ell\weak N(0,\mathbf I_r)$ in probability.  In addition, for every
continuous $f:\R^r\to\R$ of polynomial growth,
\[
    \frac1\ell\sum_{i=1}^{\ell}f(\mathbf z_i)
    \longrightarrow
    \E f(\boldsymbol{\xi}),
    \qquad
    \boldsymbol{\xi}\sim N(0,\mathbf I_r),
\]
in probability.
\label{lem:empirical-haar-row}
\end{lemma}

\begin{proof}
Let $\bG_\ell\in\R^{\ell\times r}$ have independent $N(0,1)$ entries, with rows
$\mathbf g_i$.  The Gaussian representation of the Haar Stiefel matrix gives
\[
    \bU_\ell
    \stackrel{d}{=}
    \bG_\ell(\bG_\ell^\top\bG_\ell)^{-1/2}.
\]
Hence
\[
    \mathbf z_i
    \stackrel{d}{=}
    \mathbf g_i
    \left(\frac{\bG_\ell^\top\bG_\ell}{\ell}\right)^{-1/2}.
\]
Since $r$ is fixed,
$\ell^{-1}\bG_\ell^\top\bG_\ell\to\mathbf I_r$ almost surely.  The empirical
law of the independent rows $\mathbf g_i$ converges weakly to
$N(0,\mathbf I_r)$ almost surely, and the common multiplier above converges to
the identity.  More explicitly, if
$\mathbf R_\ell=(\ell^{-1}\bG_\ell^\top\bG_\ell)^{-1/2}$, then
$\|\mathbf R_\ell-\mathbf I_r\|_{\op}\to0$ almost surely and
\[
    \frac1\ell\sum_{i=1}^{\ell}
    \varphi(\mathbf g_i\mathbf R_\ell)
    -
    \frac1\ell\sum_{i=1}^{\ell}
    \varphi(\mathbf g_i)
    \longrightarrow0
\]
for every bounded Lipschitz $\varphi$, because the Gaussian rows have bounded
empirical first moment with probability tending to one.  This proves weak
convergence of $\mu_\ell$.  For functions of polynomial growth, truncate $f$
outside a large ball and use the convergence of Gaussian moments of fixed order
to remove the truncation.
\end{proof}

\begin{lemma}[Haar-row extreme values]
Let $r$ be fixed and let
$\bU_\ell=[\bu_1,\ldots,\bu_r]\sim\mathrm{Haar}(\St(\ell,r))$.  Put
\[
    \mathbf z_i
    :=
    \sqrt\ell\,(u_1(i),\ldots,u_r(i))\in\R^r.
\]
If $\bB_\ell$ is a symmetric $r\times r$ random matrix, independent of
$\bU_\ell$, and
$\rho_\ell:=\|\bB_\ell\|_{\op}>0$ almost surely, then
\begin{equation}
    \frac{1}{2\rho_\ell\log\ell}
    \max_{1\leq i\leq\ell}
    |\mathbf z_i^\top\bB_\ell\mathbf z_i|
    \longrightarrow 1
    \qquad\text{in probability.}
    \label{eq:haar-quadratic-maximum}
\end{equation}
In particular,
\begin{equation}
    \max_{1\leq i\leq\ell}\|\mathbf z_i\|_2^2=O_{\Prob}(\log\ell),
    \qquad
    \max_{1\leq i\leq\ell}\ell u_1(i)^2=O_{\Prob}(\log\ell).
    \label{eq:haar-coordinate-maximum}
\end{equation}
\label{lem:haar-row-extremes}
\end{lemma}

\begin{proof}
It is enough to prove the quadratic-form statement uniformly over deterministic
symmetric matrices $\bB$ with $\|\bB\|_{\op}=1$; conditioning on
$\bB_\ell/\rho_\ell$ then gives the random-matrix version.  Use the Gaussian
representation
\[
    \bU_\ell
    \stackrel{d}{=}
    \mathbf G_\ell(\mathbf G_\ell^\top\mathbf G_\ell)^{-1/2},
    \qquad
    \mathbf z_i
    \stackrel{d}{=}
    \mathbf g_i\mathbf R_\ell,
    \qquad
    \mathbf R_\ell
    =
    \left(\frac{\mathbf G_\ell^\top\mathbf G_\ell}{\ell}\right)^{-1/2},
\]
where the rows $\mathbf g_i$ are independent $N(0,\mathbf I_r)$.  Since $r$ is
fixed, $\|\mathbf R_\ell-\mathbf I_r\|_{\op}\to0$ in probability.  Also
\[
    \max_i\|\mathbf g_i\|_2^2=O_{\Prob}(\log\ell)
\]
by a union bound and the chi-square tail.  Hence, uniformly over
$\|\bB\|_{\op}=1$,
\[
    \max_i
    \left|
    \mathbf g_i^\top\mathbf R_\ell\bB\mathbf R_\ell\mathbf g_i
    -
    \mathbf g_i^\top\bB\mathbf g_i
    \right|
    =
    o_{\Prob}(\log\ell).
    \label{eq:haar-extreme-perturbation}
\]

It remains to recall the standard extreme-value estimate for the independent
Gaussian quadratic forms
$Q_i(\bB):=\mathbf g_i^\top\bB\mathbf g_i$.  Diagonalize $\bB$ and let
$\rho_+(\bB)$ and $\rho_-(\bB)$ be the largest positive and negative absolute
eigenvalue magnitudes.  The tail of $Q_i(\bB)$ is governed by the largest
positive eigenspace, and the tail of $-Q_i(\bB)$ by the largest negative
eigenspace.  Since $\max\{\rho_+(\bB),\rho_-(\bB)\}=1$, the chi-square tail
asymptotic gives, uniformly over $\|\bB\|_{\op}=1$,
\[
    \max_{1\leq i\leq\ell}|Q_i(\bB)|
    =
    (2+o_{\Prob}(1))\log\ell.
\]
For completeness, the upper bound follows from
$|Q_i(\bB)|\leq\|\mathbf g_i\|_2^2$.  For the lower bound, choose an eigenvector
with eigenvalue of magnitude one.  On the event that the corresponding squared
Gaussian coordinate is at least $a\log\ell+M$ while the squared norm in the
orthogonal complement is at most $M$, one has
$|Q_i(\bB)|\geq a\log\ell$.  This event has probability at least
$c_M\ell^{-a/2}/\sqrt{\log\ell}$ for every fixed $a<2$ and some $c_M>0$ after
choosing $M$ fixed, so at least one row satisfies it with probability tending
to one.  Combining this lower bound for every $a<2$, the chi-square upper bound
for every $a>2$, and \eqref{eq:haar-extreme-perturbation} proves
\eqref{eq:haar-quadratic-maximum}.

Taking $\bB=\mathbf I_r$ gives the first bound in
\eqref{eq:haar-coordinate-maximum}, and taking $\bB=\be_1\be_1^\top$ gives the
second.
\end{proof}

\paragraph{Exact diagonal formulas.}
For each token $i$, define the left singular loading vectors
\begin{equation}
    \bar{\bu}_{Q,i}
    :=
    \bSigma_Q\bU_Q^\top\be_i\in\R^{r_Q},
    \qquad
    \bar{\bu}_{K,i}
    :=
    \bSigma_K\bU_K^\top\be_i\in\R^{r_K}.
    \label{eq:bar-loadings}
\end{equation}
These vectors include both singular values and token-side leverage.  With the
overlap matrix $\bH$ from \eqref{eq:H}, every entry of $\bQ\bK^\top$ has the form
\begin{equation}
    (\bQ\bK^\top)_{ij}
    =
    \bar{\bu}_{Q,i}^{\top}\bH\,\bar{\bu}_{K,j}.
    \label{eq:entry-form}
\end{equation}

\begin{proposition}[Exact eigenvalues for lower-triangular masks]
Assume $\bL$ is lower triangular.  Then the eigenvalues of
$\bM=\bL\hadamard\bQ\bK^\top$ are exactly
\begin{equation}
    \lambda_i(\bM)
    =
    L_{ii}\bar{\bu}_{Q,i}^{\top}\bH\,\bar{\bu}_{K,i},
    \qquad i=1,\ldots,\ell.
    \label{eq:eigenvalues-exact}
\end{equation}
Equivalently,
\begin{equation}
    \lambda_i(\bM)
    =
    L_{ii}
    \sum_{a=1}^{r_Q}\sum_{b=1}^{r_K}
    \sigma^Q_a\sigma^K_b\,
    u^Q_a(i)u^K_b(i)\,
    \langle\bv^Q_a,\bv^K_b\rangle.
    \label{eq:eigenvalue-expansion}
\end{equation}
If Assumption~\ref{ass:haar} also holds and $r_Q,r_K$ are fixed as
$d\to\infty$, then with
\[
    Z_{ab}^{(d)}
    :=
    \sqrt d\,\langle\bv^Q_a,\bv^K_b\rangle,
    \qquad
    c_{ab}^{(d)}(i)
    :=
    \frac{L_{ii}\sigma^Q_a\sigma^K_bu^Q_a(i)u^K_b(i)}{\sqrt d},
\]
one has the exact representation
\begin{equation}
    \lambda_i(\bM)
    =
    \sum_{a=1}^{r_Q}\sum_{b=1}^{r_K}
    c_{ab}^{(d)}(i)Z_{ab}^{(d)}.
    \label{eq:fixed-rank-exact-characterization}
\end{equation}
Moreover,
\[
    (Z_{ab}^{(d)})_{a,b}\weak (G_{ab})_{a,b},
\]
where the entries $G_{ab}$ are independent $N(0,1)$.
\label{prop:exact-lower-triangular-eigenvalues}
\end{proposition}

\begin{proof}
Since $\bL$ is lower triangular, $\bM=\bL\hadamard\bQ\bK^\top$ is lower
triangular.  The eigenvalues of a triangular matrix are its diagonal entries,
which gives \eqref{eq:eigenvalues-exact}.  Expanding \eqref{eq:svd} gives
\[
    \bQ\bK^\top
    =
    \sum_{a=1}^{r_Q}\sum_{b=1}^{r_K}
    \sigma^Q_a\sigma^K_b
    \langle\bv^Q_a,\bv^K_b\rangle
    \bu^Q_a(\bu^K_b)^\top,
\]
and taking the $(i,i)$ entry gives \eqref{eq:eigenvalue-expansion}.  Also,
\[
    \bar{\bu}_{Q,i}^{\top}\bH\,\bar{\bu}_{K,i}
    =
    \sum_{a=1}^{r_Q}\sum_{b=1}^{r_K}
    \sigma^Q_a\sigma^K_bu^Q_a(i)u^K_b(i)
    \langle\bv^Q_a,\bv^K_b\rangle,
\]
which identifies the compact loading-vector form with the expanded form.  The
exact representation
\eqref{eq:fixed-rank-exact-characterization} is a rewriting of
\eqref{eq:eigenvalue-expansion}, and the convergence of
$(Z_{ab}^{(d)})_{a,b}$ follows from Lemma~\ref{lem:haar-overlap}.
\end{proof}

\begin{proof}[Proof of Theorem~\ref{thm:lower-triangular-eigenvalues}]
By Proposition~\ref{prop:exact-lower-triangular-eigenvalues} and
Assumption~\ref{ass:common-left-haar}, and using that $\bL$ is unit-diagonal,
\[
    \lambda_i(\bM)
    =
    \sum_{a=1}^{r_Q}\sum_{b=1}^{r_K}
    \frac{\sigma^Q_a\sigma^K_b}{\ell\sqrt d}
    \{\ell u_a(i)u_b(i)\}
    Z_{ab}^{(d)}.
\]
The factors in this display have three different origins.  The deterministic
singular-value factor converges by assumption:
\[
    \frac{\sigma^Q_a\sigma^K_b}{\ell\sqrt d}
    \longrightarrow
    \theta_{ab}.
\]
The token factor is shared across $a,b$ at a fixed token but varies with $i$.
For any fixed collection of distinct tokens $i_1,\ldots,i_m$,
Lemma~\ref{lem:token-haar-products} gives
\[
    \left(\ell u_a(i_p)u_b(i_p)\right)_{1\leq p\leq m,\;a,b}
    \weak
    \left(\xi_{i_p,a}\xi_{i_p,b}\right)_{1\leq p\leq m,\;a,b}.
\]
The right-space factor is common to all tokens:
\[
    Z_{ab}^{(d)}
    =
    \sqrt d\,\langle\bv^Q_a,\bv^K_b\rangle,
    \qquad
    (Z_{ab}^{(d)})_{a,b}\weak (G_{ab})_{a,b}
\]
by Lemma~\ref{lem:haar-overlap}.  Assumptions~\ref{ass:haar}
and~\ref{ass:common-left-haar} make the token-side Haar matrix independent of the
right singular spaces, so the two convergences are joint with independent
limits.  Applying the continuous mapping theorem to the finite polynomial map
that sends these arrays to
$(\lambda_{i_1},\ldots,\lambda_{i_m})$ gives
\eqref{eq:common-haar-eigenvalue-limit}.  The same Gaussian matrix
$(G_{ab})$ appears for every token because the right singular spaces are shared
by all rows of $\bQ\bK^\top$.

It remains to compute the covariance in the limit.  Write
\[
    S_i
    :=
    \sum_{a=1}^{r_Q}\sum_{b=1}^{r_K}
    \theta_{ab}\xi_{i,a}\xi_{i,b}G_{ab},
    \qquad
    \Lambda_i=S_i.
\]
Since $\E G_{ab}=0$, one has $\E S_i=0$.  For distinct tokens $i_p\neq i_q$,
independence of the token-side vectors gives
\[
    \E S_{i_p}S_{i_q}
    =
    \sum_{a,b,c,d}
    \theta_{ab}\theta_{cd}
    \E[G_{ab}G_{cd}]
    \E[\xi_{i_p,a}\xi_{i_p,b}]
    \E[\xi_{i_q,c}\xi_{i_q,d}].
\]
Using
$\E[G_{ab}G_{cd}]=\one_{\{a=c,\;b=d\}}$ and
$\E[\xi_{i,a}\xi_{i,b}]=\one_{\{a=b\}}$, only the diagonal terms $a=b$ survive.
Thus
\[
    \E S_{i_p}S_{i_q}
    =
    \sum_{a=1}^{\min(r_Q,r_K)}\theta_{aa}^2
    =
    D_{\boldsymbol{\theta}},
    \qquad p\neq q.
\]
For a single token,
\[
    \E S_i^2
    =
    \sum_{a=1}^{r_Q}\sum_{b=1}^{r_K}
    \theta_{ab}^2\E[\xi_{i,a}^2\xi_{i,b}^2]
    =
    F_{\boldsymbol{\theta}}+2D_{\boldsymbol{\theta}},
\]
because $\E\xi^4=3$ on the diagonal $a=b$ and
$\E[\xi_a^2\xi_b^2]=1$ off the diagonal.  This gives
\eqref{eq:common-haar-eigenvalue-covariance}.
\end{proof}

\subsection{Eigenvalue Distribution}
\label{lr:eigenvalue-details}

For the finite model on the joint sequence, embed the coefficient matrix in
$\R^{r\times r}$ by setting
\begin{equation}
    A_{ab}^{(n)}
    :=
    \theta_{ab}^{(n)}Z_{ab}^{(d_n)},
    \qquad
    1\leq a\leq r_Q,\quad 1\leq b\leq r_K,
    \qquad
    \bB_n:=\frac{\bA_n+\bA_n^\top}{2},
    \label{eq:finite-eigenvalue-coefficient-matrix}
\end{equation}
and setting every other entry of $\bA_n$ to zero.  With
$\mathbf z_i=\sqrt{\ell_n}\,(u_1(i),\ldots,u_r(i))$, the exact diagonal formula
becomes
\begin{equation}
    \lambda_i(\bM)=\mathbf z_i^\top\bB_n\mathbf z_i.
    \label{eq:finite-eigenvalue-quadratic-form}
\end{equation}
The deterministic singular-value convention and the frame assumptions make
$\bB_n$ independent of $\bU$.  Moreover,
\begin{equation}
    \bB_n\weak\bB(\bG)
    :=\frac{\bA(\bG)+\bA(\bG)^\top}{2},
    \qquad
    A_{ab}(\bG)=\theta_{ab}G_{ab},
    \label{eq:finite-coefficient-matrix-limit}
\end{equation}
where $\bA(\bG)\in\R^{r\times r}$ is likewise zero outside its
$r_Q\times r_K$ block.  For a symmetric matrix $\bC\in\R^{r\times r}$, write
\[
    \nu_{\bC}:=\mathcal L_{\boldsymbol{\xi}}
    (\boldsymbol{\xi}^\top\bC\boldsymbol{\xi}),
    \qquad \boldsymbol{\xi}\sim N(0,\mathbf I_r).
\]

\begin{proposition}[Non-collapse of the eigenvalue empirical law]
Assume the setting of Theorem~\ref{thm:lower-triangular-eigenvalues},
$\boldsymbol{\theta}\neq0$.  Let
$\mu_{\mathrm{eig},\ell}:=\ell^{-1}\sum_{i=1}^{\ell}
\delta_{\lambda_i(\bM)}$, and let $d_{\mathrm{BL}}$ denote the bounded-Lipschitz
metric on probability measures.  Then
\begin{equation}
    d_{\mathrm{BL}}\!\left(\mu_{\mathrm{eig},\ell},\nu_{\bB_n}\right)
    \longrightarrow 0
    \quad\text{in probability},
    \qquad
    (\bB_n,\mu_{\mathrm{eig},\ell})
    \weak
    (\bB(\bG),\nu_{\bB(\bG)})
    \label{eq:noncollapsed-esd-limit}
\end{equation}
as a matrix--random-measure pair.  Moreover
$\nu_{\bB(\bG)}\neq\delta_0$ almost surely, so
$\mu_{\mathrm{eig},\ell}$ does not converge weakly to $\delta_0$ in
probability.
\label{prop:eigenvalue-esd-noncollapse}
\end{proposition}

\begin{proof}
Condition on the finite coefficient matrix $\bB_n$.  By
Lemma~\ref{lem:empirical-haar-row}, the empirical law of the quadratic forms in
\eqref{eq:finite-eigenvalue-quadratic-form} approaches $\nu_{\bB_n}$.  The
convergence is uniform with $\bB_n$ in compact sets; since $\bB_n$ is tight,
this proves the first assertion in \eqref{eq:noncollapsed-esd-limit}.  The map
$\bC\mapsto\nu_{\bC}$ is continuous for the weak topology, so
\eqref{eq:finite-coefficient-matrix-limit} and Slutsky's theorem give the joint
convergence in \eqref{eq:noncollapsed-esd-limit}.  If
$\boldsymbol{\theta}\neq0$, the zero-padded symmetric Gaussian matrix
$\bB(\bG)$ is nonzero almost surely.  Its Gaussian quadratic form is therefore
not identically zero, and $\nu_{\bB(\bG)}\neq\delta_0$ almost surely.
\end{proof}

\begin{proposition}[Effective eigenvalue counts and tail probabilities]
Assume the setting of Theorem~\ref{thm:lower-triangular-eigenvalues},
$\boldsymbol{\theta}\neq0$.  Define
\[
    \barPhi(t):=\Prob\{Z>t\},
    \qquad Z\sim N(0,1),
\]
and
\[
    N_\varepsilon^{|\cdot|}
    :=
    \#\{i:|\lambda_i(\bM)|>\varepsilon\}.
\]
For $\boldsymbol{\xi}\sim N(0,\mathbf I_r)$, set
\[
    V_{\boldsymbol{\theta}}(\boldsymbol{\xi})
    :=
    \sum_{a=1}^{r_Q}\sum_{b=1}^{r_K}
    \theta_{ab}^2\xi_a^2\xi_b^2,
    \qquad
    s_{\boldsymbol{\theta}}^2
    :=
    \E V_{\boldsymbol{\theta}}(\boldsymbol{\xi})
    =
    \sum_{a,b}(1+2\mathbf 1_{\{a=b\}})\theta_{ab}^2.
\]
For $\varepsilon>0$, define
\begin{equation}
    p_\varepsilon^+(\boldsymbol{\theta})
    =
    \E_{\boldsymbol{\xi}}
    \barPhi\left(
        \frac{\varepsilon}
        {\sqrt{V_{\boldsymbol{\theta}}(\boldsymbol{\xi})}}
    \right),
    \qquad
    p_\varepsilon^{|\cdot|}(\boldsymbol{\theta})
    =
    2p_\varepsilon^+(\boldsymbol{\theta}),
    \label{eq:eigenvalue-tail-definitions}
\end{equation}
where the integrand is set equal to zero when
$V_{\boldsymbol{\theta}}(\boldsymbol{\xi})=0$.  Then, for every
$\varepsilon>0$,
\begin{equation}
    \frac{\E N_\varepsilon^{|\cdot|}}{\ell}
    \longrightarrow
    p_\varepsilon^{|\cdot|}(\boldsymbol{\theta}).
    \label{eq:effective-count-exact-limit}
\end{equation}
Moreover, for every $u>0$,
\begin{equation}
    \frac{3}{70}
    \barPhi\left(
    \frac{\sqrt2\,\varepsilon}{s_{\boldsymbol{\theta}}}
    \right)
    \leq
    p_\varepsilon^{|\cdot|}(\boldsymbol{\theta})
    \leq
    2\barPhi(\varepsilon/u)+2s_{\boldsymbol{\theta}}^2/u^2.
    \label{eq:eigenvalue-tail-two-sided}
\end{equation}
Also,
\begin{equation}
    p_\varepsilon^{|\cdot|}(\boldsymbol{\theta})
    \leq
    C\exp\left\{
        -c\left(\frac{\varepsilon}{s_{\boldsymbol{\theta}}}\right)^{2/3}
    \right\},
    \label{eq:eigenvalue-chaos-tail-bound}
\end{equation}
for universal constants $c,C>0$.  For the finite coefficient matrix, set
\[
    q_{\varepsilon,n}
    :=
    \Prob_{\boldsymbol{\xi}}
    \left\{\left|\boldsymbol{\xi}^\top\bB_n\boldsymbol{\xi}\right|>
    \varepsilon\right\}.
\]
Then
\begin{equation}
    \frac{N_\varepsilon^{|\cdot|}}{\ell}-q_{\varepsilon,n}
    \longrightarrow 0
    \quad\text{in probability},
    \qquad
    \left(\bB_n,\frac{N_\varepsilon^{|\cdot|}}{\ell}\right)
    \weak
    \left(\bB(\bG),
    p_\varepsilon^{|\cdot|}(\boldsymbol{\theta};\bG)\right),
    \label{eq:conditional-effective-count}
\end{equation}
where
\[
    p_\varepsilon^{|\cdot|}(\boldsymbol{\theta};\bG)
    :=
    \Prob_{\boldsymbol{\xi}}
    \left\{
    \left|
    \sum_{a,b}\theta_{ab}\xi_a\xi_bG_{ab}
    \right|>\varepsilon
    \right\}.
\]
In particular,
$N_\varepsilon^{|\cdot|}/\ell\weak
p_\varepsilon^{|\cdot|}(\boldsymbol{\theta};\bG)$.  In the rank-one case,
\begin{equation}
    p_\varepsilon^+(\theta)
    =
    \E_{\xi}
    \barPhi\left(\frac{\varepsilon}{|\theta|\xi^2}\right)
    =
    \int_0^\infty
    \barPhi\left(\frac{\varepsilon}{|\theta|x}\right)
    \frac{e^{-x/2}}{\sqrt{2\pi x}}\,dx.
    \label{eq:rank-one-tail-formula}
\end{equation}
\label{prop:eigenvalue-effective-counts}
\end{proposition}

\begin{proof}
Theorem~\ref{thm:lower-triangular-eigenvalues} gives, for a fixed token,
\[
    \lambda_i(\bM)
    \weak
    \sum_{a=1}^{r_Q}\sum_{b=1}^{r_K}
    \theta_{ab}\xi_a\xi_bG_{ab},
\]
where $(G_{ab})$ is independent of $\boldsymbol{\xi}$.  Conditional on
$\boldsymbol{\xi}$, the right-hand side is a centered Gaussian random variable
with variance $V_{\boldsymbol{\theta}}(\boldsymbol{\xi})$.  Hence
\[
    \Prob_{\bG}
    \left\{
    \sum_{a,b}\theta_{ab}\xi_a\xi_bG_{ab}>\varepsilon
    \,\middle|\,\boldsymbol{\xi}
    \right\}
    =
    \barPhi\left(
        \frac{\varepsilon}{\sqrt{V_{\boldsymbol{\theta}}(\boldsymbol{\xi})}}
    \right),
\]
with value zero when the variance is zero.  Averaging over
$\boldsymbol{\xi}$ gives the one-sided formula defining
$p_\varepsilon^+(\boldsymbol{\theta})$.  The limiting variable is symmetric
about zero after averaging over $\bG$, so the two-sided formula is twice the
one-sided formula.

Linearity of expectation and exchangeability of the Haar token rows give
\[
    \frac{\E N_\varepsilon^{|\cdot|}}{\ell}
    =
    \Prob\{|\lambda_1(\bM)|>\varepsilon\}
    \longrightarrow
    p_\varepsilon^{|\cdot|}(\boldsymbol{\theta}),
\]
which proves \eqref{eq:effective-count-exact-limit}.  The indicator
discontinuity is harmless: for $\varepsilon>0$ the limiting Gaussian-polynomial
variable has no atom at $\pm\varepsilon$ unless it is identically zero, and the
latter has probability zero when $\boldsymbol{\theta}\neq0$.

For the upper estimate in \eqref{eq:eigenvalue-tail-two-sided}, split on
$V_{\boldsymbol{\theta}}(\boldsymbol{\xi})\leq u^2$.  On this event the
integrand is at most $\barPhi(\varepsilon/u)$, while the complement has
probability at most $s_{\boldsymbol{\theta}}^2/u^2$ by Markov's inequality.

For the lower estimate, Paley--Zygmund applied to
$V_{\boldsymbol{\theta}}(\boldsymbol{\xi})$ gives, with
$K_{\boldsymbol{\theta}}:=
\E\{V_{\boldsymbol{\theta}}(\boldsymbol{\xi})^2\}/s_{\boldsymbol{\theta}}^4$,
\[
    \Prob\left\{
    V_{\boldsymbol{\theta}}(\boldsymbol{\xi})
    \geq
    \frac{s_{\boldsymbol{\theta}}^2}{2}
    \right\}
    \geq
    \frac{1}{4K_{\boldsymbol{\theta}}}.
\]
On this event,
$\barPhi\{\varepsilon/V_{\boldsymbol{\theta}}(\boldsymbol{\xi})^{1/2}\}
\geq
\barPhi(\sqrt2\,\varepsilon/s_{\boldsymbol{\theta}})$.  Multiplying by the
outer factor two in $p_\varepsilon^{|\cdot|}=2p_\varepsilon^+$ gives the lower
bound in \eqref{eq:eigenvalue-tail-two-sided}.  Finally,
$K_{\boldsymbol{\theta}}\leq 35/3$ because
\[
    \frac{
    \E\{\xi_a^2\xi_b^2\xi_{a'}^2\xi_{b'}^2\}}
    {(1+2\mathbf 1_{\{a=b\}})(1+2\mathbf 1_{\{a'=b'\}})}
    \leq
    \frac{35}{3}
\]
for all indices; the maximum occurs when $a=b=a'=b'$.

The Gaussian chaos bound \eqref{eq:eigenvalue-chaos-tail-bound} is the standard
hypercontractive tail estimate for a Gaussian polynomial of degree at most
three, applied to $\sum_{a,b}\theta_{ab}\xi_a\xi_bG_{ab}$.  For the empirical
count, condition on the finite matrix $\bB_n$ and apply
Lemma~\ref{lem:empirical-haar-row} to the token variables.  Indicator functions
are handled by approximation at continuity points.  Every nonzero Gaussian
quadratic form has no atom at $\pm\varepsilon$, and $\bB_n\neq0$ almost surely.
Thus $N_\varepsilon^{|\cdot|}/\ell-q_{\varepsilon,n}\to0$ in probability.
Finally, the map
$\bC\mapsto\Prob_{\boldsymbol{\xi}}\{|\boldsymbol{\xi}^\top\bC
\boldsymbol{\xi}|>\varepsilon\}$ is continuous at every nonzero $\bC$.
Combining this fact with \eqref{eq:finite-coefficient-matrix-limit} proves the
joint convergence in \eqref{eq:conditional-effective-count}.
In the rank-one case, write $x=\xi^2$ and use the $\chi^2(1)$ density to obtain
\eqref{eq:rank-one-tail-formula}.
\end{proof}

\begin{proof}[Proof of Theorem~\ref{thm:eigenvalue-spectral-phenomena}]
Part (i) follows from Proposition~\ref{prop:eigenvalue-effective-counts} and the
lower bound \eqref{eq:eigenvalue-tail-two-sided}.  It remains to prove the rank
statements in part (ii).

Use the finite coefficient matrix $\bB_n$ from
\eqref{eq:finite-eigenvalue-coefficient-matrix}, and put
$\rho_n:=\|\bB_n\|_{\op}$.  For any nonzero symmetric matrix $\bC$, define
\[
    \rho(\bC):=\|\bC\|_{\op},
    \qquad
    Y_{\bC}:=\boldsymbol{\xi}^\top\bC\boldsymbol{\xi},
    \qquad
    \mu_1(\bC):=\E_{\boldsymbol{\xi}}|Y_{\bC}|,
    \qquad
    \eta_1(\bC):=
    \E_{\boldsymbol{\xi}}\{|Y_{\bC}|\log|Y_{\bC}|\},
\]
with $0\log0=0$.  Write $\rho_*(\bG):=\rho(\bB(\bG))$,
$Y_{\bG}:=Y_{\bB(\bG)}$, $\mu_1(\bG):=\mu_1(\bB(\bG))$, and
$\eta_1(\bG):=\eta_1(\bB(\bG))$.  Since
$\boldsymbol{\theta}\neq0$, both $\bB_n\neq0$ almost surely and
$\bB(\bG)\neq0$ almost surely.  Lemma~\ref{lem:haar-row-extremes} gives
\[
    \frac{\max_i|\lambda_i(\bM)|}{2\rho_n\log\ell}
    \longrightarrow 1
    \qquad\text{in probability.}
\]
Conditioning on $\bB_n$ and applying
Lemma~\ref{lem:empirical-haar-row} to the continuous polynomial-growth
functions $|\mathbf z^\top\bB_n\mathbf z|$ and
$|\mathbf z^\top\bB_n\mathbf z|\log|\mathbf z^\top\bB_n\mathbf z|$ gives
\begin{align}
    \frac{1}{\ell}\sum_{i=1}^{\ell}|\lambda_i(\bM)|
    -\mu_1(\bB_n)&\longrightarrow0,
    \label{eq:finite-eigenvalue-first-moment}\\
    \frac{1}{\ell}\sum_{i=1}^{\ell}
    |\lambda_i(\bM)|\log|\lambda_i(\bM)|
    -\eta_1(\bB_n)&\longrightarrow0
    \label{eq:finite-eigenvalue-log-moment}
\end{align}
in probability.  As in the empirical-law proof, the conditional convergence is
uniform on compact sets of coefficient matrices, and $\bB_n$ is tight.  These
convergences therefore hold jointly with
$\bB_n\weak\bB(\bG)$.  Consequently,
\begin{equation}
    \frac{\log\ell}{\ell}\srank_{\lambda}(\bM)
    \weak
    \frac{\mu_1(\bG)}{2\rho_*(\bG)}.
    \label{eq:eigenvalue-stable-rank-prefactor}
\end{equation}

Diagonalizing $\bB(\bG)$ gives
\[
    Y_{\bG}
    \stackrel{d}{=}
    \sum_{k=1}^r\rho_k(\bG)\chi_k^2,
\]
where $\rho_k(\bG)$ are the eigenvalues of $\bB(\bG)$ and
$\rho_*(\bG)=\max_k|\rho_k(\bG)|$.
The upper bound on the prefactor follows from
\[
    \mu_1(\bG)
    \leq
    \E_{\boldsymbol{\xi}}\sum_{k=1}^r|\rho_k(\bG)|\chi_k^2
    =
    \sum_{k=1}^r|\rho_k(\bG)|
    \leq
    r\rho_*(\bG).
\]
This gives the displayed upper bound after dividing by $2\rho_*(\bG)$.  For
the lower bound, let $Y=\boldsymbol{\xi}^\top\bC\boldsymbol{\xi}$ for a
nonzero symmetric $\bC$.  Gaussian hypercontractivity and log-convexity of
$L^p$ norms give
\[
    \|Y\|_4\leq3\|Y\|_2,
    \qquad
    \|Y\|_2\leq\|Y\|_1^{1/3}\|Y\|_4^{2/3},
\]
and hence $\E|Y|\geq\|Y\|_2/9$.  Moreover,
\[
    \E Y^2
    =2\operatorname{tr}(\bC^2)+(\operatorname{tr}\bC)^2
    \geq2\|\bC\|_{\op}^2.
\]
Taking $\bC=\bB(\bG)$ yields the claimed universal bound
$\mu_1(\bG)\geq \sqrt2\,\rho_*(\bG)/9$.  This proves
the $\Theta_{\Prob}(\ell/\log\ell)$ statement because $\bB(\bG)\neq0$ almost
surely when $\boldsymbol{\theta}\neq0$.

Since
\[
    -\sum_{i=1}^{\ell}p_i\log p_i
    =
    \log\left(\sum_{i=1}^{\ell}w_i\right)
    -
    \frac{\sum_{i=1}^{\ell}w_i\log w_i}
    {\sum_{i=1}^{\ell}w_i},
\]
we obtain
\begin{align}
    \log\operatorname{erank}_{\lambda}(\bM)-\log\ell
    &=
    \log\mu_1(\bB_n)
    -\frac{\eta_1(\bB_n)}{\mu_1(\bB_n)}
    +o_{\Prob}(1),
    \notag\\
    \log\operatorname{erank}_{\lambda}(\bM)-\log\ell
    &\weak
    \log\mu_1(\bG)-\frac{\eta_1(\bG)}{\mu_1(\bG)},
    \label{eq:eigenvalue-entropy-limit}
\end{align}
where the second conclusion follows from
\eqref{eq:finite-coefficient-matrix-limit} by the continuous mapping theorem.
Finally, put $W=|Y_{\bG}|$ and $X=W/\mu_1(\bG)$.  Jensen's inequality under
the probability measure with density $X$ gives
\[
    \frac{\eta_1(\bG)}{\mu_1(\bG)}-\log\mu_1(\bG)
    =\E(X\log X)
    \leq\log\E X^2
    \leq\log81,
\]
where the last inequality follows from
$\mu_1(\bG)\geq\|Y_{\bG}\|_2/9$.  Thus
\[
    \log\mu_1(\bG)
    -
    \frac{\eta_1(\bG)}{\mu_1(\bG)}
    \geq
    -\log81,
\]
and the entropy of a probability vector on $\ell$ points is at most $\log\ell$.
This proves the linear lower and upper bounds in
\eqref{eq:eigenvalue-rank-summary}.
\end{proof}

\begin{remark}[Signal extraction threshold comparison]
For the structured causal masks above, the common-Haar token calculation gives
the nontrivial eigenvalue-level threshold
$\sigma^Q_a\sigma^K_b\asymp\ell\sqrt d$; divergence beyond the limiting scale
requires $\sigma^Q_a\sigma^K_b\gg\ell\sqrt d$.  This is much stronger than the
softmax-attention threshold in \citet{li:2026spectral}.  In that setting,
isotropic score geometry can preserve signals of order larger than $\log \ell$.
\label{rem:threshold-comparison}
\end{remark}

\begin{remark}[Why $\rho_*$ is not one]
The quantity $\rho_*(\bG)$ above is not equal to one in this fixed-rank setting.
It is the largest absolute eigenvalue of the finite symmetric variance-profile
matrix $\bB(\bG)=\{\bA(\bG)+\bA(\bG)^\top\}/2$, where
$A_{ab}(\bG)=\theta_{ab}G_{ab}$.  The circular law concerns the eigenvalues of
large non-Hermitian Ginibre matrices after a $1/\sqrt r$ normalization.  It does
not apply here.
\label{rem:rho-star-not-one}
\end{remark}

\subsection{Singular-Value Details}
\label{lr:singular-details}

\begin{proposition}[Low-rank singular-value reduction]
Under the decomposition \eqref{eq:svd},
\begin{equation}
    \bM
    =
    \sum_{a=1}^{r_Q}
    \sum_{b=1}^{r_K}
    \sigma^Q_a\sigma^K_b
    \langle \bv^Q_a,\bv^K_b\rangle
    \diag(\bu^Q_a)\,\bL\,\diag(\bu^K_b).
    \label{eq:singular-reduction}
\end{equation}
In the rank-one case,
\[
    s_j(\bM)
    =
    |\sigma^Q\sigma^K\langle \bv^Q,\bv^K\rangle|\,
    s_j\{\diag(\bu^Q)\bL\diag(\bu^K)\}.
\]
\label{prop:singular-reduction}
\end{proposition}

\begin{proof}
From the singular value decompositions,
\[
    \bQ\bK^\top
    =
    \sum_{a=1}^{r_Q}
    \sum_{b=1}^{r_K}
    \sigma^Q_a\sigma^K_b
    \langle \bv^Q_a,\bv^K_b\rangle
    \bu^Q_a(\bu^K_b)^\top
\]
because $(\bv^Q_a)^\top\bv^K_b=\langle\bv^Q_a,\bv^K_b\rangle$.  For any pair
of vectors $\mathbf x,\mathbf y\in\R^\ell$, the Hadamard product with a rank-one
outer product can be written entrywise as
\[
    \{\bL\hadamard(\mathbf x\mathbf y^\top)\}_{ij}
    =
    L_{ij}x_i y_j
    =
    \{\diag(\mathbf x)\bL\diag(\mathbf y)\}_{ij}.
\]
Applying this identity with
$\mathbf x=\bu^Q_a$ and $\mathbf y=\bu^K_b$ gives
\[
    \bL\hadamard(\bu^Q_a(\bu^K_b)^\top)
    =
    \diag(\bu^Q_a)\bL\diag(\bu^K_b)
\]
term by term, proving \eqref{eq:singular-reduction}.  In the rank-one case the
sum has a single term, so $\bM$ is a scalar multiple of
$\diag(\bu^Q)\bL\diag(\bu^K)$.  Singular values are homogeneous under scalar
multiplication, which gives the displayed rank-one formula.
\end{proof}

\begin{proposition}[Stable semiseparable norm bound]
Let $\bL$ be the semiseparable mask \eqref{eq:semiseparable-mask} and assume
$0\leq a_t\leq a<1$ for all $t$.  Then
$\|\bL\|_{\op}\leq(1-a)^{-1}$.
\label{prop:semiseparable-bound}
\end{proposition}

\begin{proof}
For $i\geq j$, $|L_{ij}|\leq a^{i-j}$.  Hence every row sum and every column sum
of $|\bL|$ is bounded by $\sum_{k\geq0}a^k=(1-a)^{-1}$.  Schur's test gives
$\|\bL\|_{\op}\leq(1-a)^{-1}$.
\end{proof}

\begin{proposition}[Principal inverse bounds for causal masks]
Let $\bL_S$ denote the principal submatrix of $\bL$ indexed by
$\varnothing\neq S\subset\{1,\ldots,\ell\}$, with the inherited order.  For
the causal linear mask, $\|\bL_S^{-1}\|_{\op}\leq2$.  For the decay mask
$L_{ij}=\one_{\{i\geq j\}}\gamma^{i-j}$ with $0\leq\gamma<1$,
$\|\bL_S^{-1}\|_{\op}\leq1+\gamma$.  For the semiseparable mask
\eqref{eq:semiseparable-mask} with $0\leq a_t\leq a<1$,
$\|\bL_S^{-1}\|_{\op}\leq1+a$.
\label{prop:principal-inverse-bound}
\end{proposition}

\begin{proof}
Write $S=\{i_1<\cdots<i_m\}$.  In the causal linear case, $\bL_S$ is the
$m\times m$ lower triangular all-one matrix.  Its inverse has one on the
diagonal, minus one on the first subdiagonal, and zero elsewhere, so its
operator norm is at most two.

For the decay mask, $(\bL_S)_{pq}=\one_{\{p\geq q\}}\gamma^{i_p-i_q}$.  Its
inverse is lower bidiagonal with diagonal entries one and subdiagonal entries
$-\gamma^{i_p-i_{p-1}}$, whose absolute values are at most $\gamma$.  The
semiseparable case is identical after replacing
$\gamma^{i_p-i_{p-1}}$ by $\prod_{t=i_{p-1}+1}^{i_p}a_t$, whose absolute value
is at most $a$.  The displayed operator norm bounds follow from the triangle
inequality.
\end{proof}

\begin{proposition}[Weighted causal-Volterra spread]
For the causal linear mask $\bL_{ij}=\one_{\{i\geq j\}}$, write the rank-one
token vector as $\bu=(u(1),\ldots,u(\ell))$ and
$\bw=\sqrt\ell(|u(1)|,\ldots,|u(\ell)|)$, where $\bu$ is uniformly
distributed on the unit sphere in $\R^\ell$.  There exist deterministic
constants $c_V,C_V>0$, independent of $\ell$, such that, with probability
tending to one,
\begin{equation}
    \|\diag(\bw)\bL\diag(\bw)\|_*
    \geq c_V\ell\log\ell,
    \qquad
    \|\diag(\bw)\bL\diag(\bw)\|_{\op}
    \leq C_V\ell.
    \label{eq:weighted-volterra-spread}
\end{equation}
\label{prop:weighted-volterra-spread}
\end{proposition}

\begin{proof}
Put $\bT_\ell=\diag(\bw)\bL\diag(\bw)$.  Since
$\sum_iw_i^2=\ell$ exactly and Lemma~\ref{lem:empirical-haar-row} gives
$\ell^{-1}\sum_iw_i^4\to3$ in probability,
\[
    \|\bT_\ell\|_F^2
    =
    \sum_{i\geq j}w_i^2w_j^2
    =
    \frac12\ell^2+\frac12\sum_iw_i^4
    =
    \left(\frac12+o_{\Prob}(1)\right)\ell^2.
\]
Consequently $\|\bT_\ell\|_{\op}\leq\|\bT_\ell\|_F\leq C_V\ell$
with probability tending to one for a deterministic $C_V$.

It remains to prove the nuclear-norm lower bound.  Fix, for example, $t=1$ and let
$S_t=\{i:w_i>t\}$ and $m=|S_t|$.  The empirical Haar-row law gives
\[
    \frac{m}{\ell}\longrightarrow
    p_t:=\Prob\{|N(0,1)|>t\}>0
\]
in probability.  With the indices in inherited order, the principal block is
\[
    (\bT_\ell)_{S_t}=\bD_{S_t}\mathbf J_m\bD_{S_t},
    \qquad
    \bD_{S_t}:=\diag((w_i)_{i\in S_t}),
    \qquad
    (\mathbf J_m)_{pq}=\one_{\{p\geq q\}},
\]
on the event $m\geq1$, whose probability tends to one.  Every diagonal entry
of $\bD_{S_t}$ is larger than $t$.  Compression is contractive for the nuclear
norm, and the ideal property of that norm gives
\[
    \|\bT_\ell\|_*
    \geq
    \|(\bT_\ell)_{S_t}\|_*
    \geq
    t^2\|\mathbf J_m\|_*.
\]
The inverse of $\mathbf J_m$ is the first-difference matrix.  Diagonalizing
$(\mathbf J_m^{-1})^\top\mathbf J_m^{-1}$ gives
\[
    s_k(\mathbf J_m)
    =
    \left\{
        2\sin\left(\frac{(2k-1)\pi}{4m+2}\right)
    \right\}^{-1},
    \qquad k=1,\ldots,m.
\]
Using $\sin x\leq x$ and
$\sum_{k=1}^m(2k-1)^{-1}\geq\frac12\log(m+1)$,
\[
    \|\mathbf J_m\|_*
    \geq
    \frac{2m+1}{\pi}
    \sum_{k=1}^m\frac{1}{2k-1}
    \geq
    c\,m\log(m+1).
\]
Since $m/\ell\to p_t>0$ in probability, this proves
$\|\bT_\ell\|_*\geq c_V\ell\log\ell$ with probability tending to one for a
deterministic $c_V>0$.
\end{proof}

\begin{proof}[Proof of Theorem~\ref{thm:singular-spectral-phenomena}]
Let $Z_d$ and $a_{\ell,d}$ be as in the theorem.  Assumption~\ref{ass:haar}
gives $Z_d\weak G\sim N(0,1)$, while $a_{\ell,d}\to\theta$.
Writing $w_i=\sqrt\ell\,|u(i)|$ and using sign matrices, the singular values of
$\diag(\bu)\bL\diag(\bu)$ are $\ell^{-1}$ times those of
$\diag(\bw)\bL\diag(\bw)$.  Proposition~\ref{prop:singular-reduction}
therefore gives
\[
    \bT_\ell:=\diag(\bw)\bL\diag(\bw),
    \qquad
    w_i=\sqrt\ell\,|u(i)|,
    \qquad
    s_j(\bM)=|a_{\ell,d}Z_d|s_j(\bT_\ell).
\]

We first prove non-collapse.  Fix $t>0$ and set
$S_t=\{i:w_i>t\}$.  Lemma~\ref{lem:empirical-haar-row}, applied to the
continuity set $\{\mathbf z:|z_1|>t\}$, gives
$|S_t|/\ell\to p_t$.  With probability tending to one $S_t$ is nonempty, and
its principal submatrix satisfies
\[
    s_{\min}\{(\bT_\ell)_{S_t}\}
    \geq
    t^2s_{\min}(\bL_{S_t})
    \geq
    \frac{t^2}{C_{\mathrm{inv}}}.
\]
By singular-value interlacing for submatrices, $\bT_\ell$ has at least
$|S_t|$ singular values larger than $t^2/C_{\mathrm{inv}}$.  Since
$Z_d\neq0$ almost surely, the displayed formula for $s_j(\bM)$ implies
\eqref{eq:singular-noncollapse}.

For the deterministic-threshold consequence, fix $g_0>0$.  For all
sufficiently large $\ell$, $|a_{\ell,d}|\geq|\theta|/2$.  Hence, on
$\{|Z_d|\geq g_0\}$, each of the $|S_t|$ singular values just found is larger
than
\[
    \frac{|\theta|g_0t^2}{2C_{\mathrm{inv}}}
    >
    \varepsilon_{t,g_0}.
\]
The token-side vector $\bu$ is independent of $Z_d$, so conditioning on
$\{|Z_d|\geq g_0\}$ does not change the convergence $|S_t|/\ell\to p_t$ in
probability.  This proves the conditional assertion in the theorem.  Finally,
$\Prob\{|Z_d|\geq g_0\}\to\Prob\{|G|\geq g_0\}$, and independence gives the
unconditional lower bound in
\eqref{eq:singular-noncollapse-deterministic}.

For the decay and stable semiseparable masks, Schur's test gives
$\|\bL\|_{\op}\leq C_L$ with $C_L=(1-\gamma)^{-1}$ for decay and
$C_L=(1-a)^{-1}$ for semiseparable masks.  Therefore
\[
    \|\bT_\ell\|_{\op}
    \leq
    C_L\max_i w_i^2.
\]
Taking the one-dimensional quadratic form in
Lemma~\ref{lem:haar-row-extremes} gives the stronger asymptotic
\[
    \frac{\max_iw_i^2}{2\log\ell}\longrightarrow1
\]
in probability.  On the other hand, since $L_{ii}=1$,
\[
    \|\bT_\ell\|_F^2
    \geq
    \sum_{i=1}^{\ell}w_i^4,
    \qquad
    \|\bT_\ell\|_*
    \geq
    |\operatorname{Tr}\bT_\ell|
    =
    \sum_{i=1}^{\ell}w_i^2.
\]
Lemma~\ref{lem:empirical-haar-row} gives
$\ell^{-1}\sum_iw_i^4\to\E N(0,1)^4=3$ and
$\ell^{-1}\sum_iw_i^2\to1$.  Thus
\[
    \srank(\bT_\ell)
    =
    \frac{\|\bT_\ell\|_F^2}{\|\bT_\ell\|_{\op}^2}
    \geq
    c\,\frac{\ell}{(\log\ell)^2},
    \qquad
    \frac{\|\bT_\ell\|_*}{\|\bT_\ell\|_{\op}}
    \geq
    c\,\frac{\ell}{\log\ell}
\]
with probability tending to one.  Since Shannon entropy dominates min-entropy,
$\erank_{1}(\bT_\ell)\geq
\|\bT_\ell\|_*/\|\bT_\ell\|_{\op}$.  Applying the same min-entropy bound to the
squared singular-value weights gives
\[
    \erank_{2}(\bT_\ell)
    \geq
    \frac{\|\bT_\ell\|_F^2}{\|\bT_\ell\|_{\op}^2}
    =
    \srank(\bT_\ell).
\]
This proves part (ii).  The scalar prefactor
$|a_{\ell,d}Z_d|$ does not change stable rank or either entropy effective rank.

For the causal linear mask, the same empirical law gives
\[
    \|\bT_\ell\|_F^2
    =
    \sum_{i\geq j}w_i^2w_j^2
    =
    \frac12\left(\sum_iw_i^2\right)^2
    +
    \frac12\sum_iw_i^4
    =
    \Theta_{\Prob}(\ell^2).
\]
The lower bound $\|\bT_\ell\|_{\op}\geq c\ell$ follows by testing the matrix on
the normalized all-one vector and using
$\ell^{-2}\sum_{i\geq j}w_iw_j\to(\E|N(0,1)|)^2/2=1/\pi$.
Proposition~\ref{prop:weighted-volterra-spread} supplies
$\|\bT_\ell\|_{\op}\leq C_V\ell$ and
$\|\bT_\ell\|_*\geq c_V\ell\log\ell$ with probability tending to one.  Hence
$\srank(\bT_\ell)=\Theta_{\Prob}(1)$, and
\[
    \erank_{1}(\bT_\ell)
    \geq
    \frac{\|\bT_\ell\|_*}{\|\bT_\ell\|_{\op}}
    \geq
    c\log\ell,
\]
and the squared-weight min-entropy bound gives
\[
    \erank_{2}(\bT_\ell)
    \geq
    \srank(\bT_\ell)
    \geq c
\]
with probability tending to one.  This proves part (iii).
\end{proof}


\section{Softmax Attention Comparisons}
\label{app:softmax-comparisons}

This appendix records the corresponding row-softmax comparisons.  These models
are not instances of the entrywise masked model \eqref{eq:model}: the softmax
denominator couples all entries in a row.  Given scores
\[
    \bS=\frac{\bQ\bK^\top}{\tau_\ell},
\]
the dense self-attention matrix is
\[
    \bP=\operatorname{softmax}_{\mathrm{row}}(\bS),
    \qquad
    P_{ij}=\frac{e^{S_{ij}}}{\sum_{m=1}^{\ell}e^{S_{im}}},
\]
while the causal self-attention matrix is
\begin{equation}
    \bP=\operatorname{softmax}_{\mathrm{row}}^{\mathrm{causal}}(\bS),
    \qquad
    P_{ij}
    =
    \frac{\one_{\{j\leq i\}}e^{S_{ij}}}
    {\sum_{m=1}^{i}e^{S_{im}}}.
    \label{eq:self-attention-softmax}
\end{equation}
The causal all-one mask selects the support of \eqref{eq:self-attention-softmax}
by excluding future tokens from the denominator, but the resulting matrix is a
normalized stochastic operator rather than $\bL\hadamard\bS$.

\subsection{Dense Low-Rank Softmax}

\begin{theorem}[Dense low-rank softmax has bounded rank proxies]
\label{thm:dense-low-rank-softmax-collapse}
Let
\[
    \bP_\ell=\operatorname{softmax}_{\mathrm{row}}(\bS_\ell),
    \qquad
    P_{ij}=\frac{e^{S_{ij}}}{\sum_{m=1}^{\ell}e^{S_{im}}}.
\]
Assume that the score matrix has fixed-rank bounded feature form
\begin{equation}
    S_{ij}=\mathbf x_i^\top\mathbf y_j,\qquad
    \mathbf x_i,\mathbf y_j\in\R^r,\qquad
    \max_i\|\mathbf x_i\|_2\vee\max_j\|\mathbf y_j\|_2\leq R,
    \label{eq:bounded-low-rank-dense-softmax-scores}
\end{equation}
where $r$ and $R$ do not depend on $\ell$.  Let
\[
    \mu_{\mathrm{eig},\ell}^{\mathrm{sm}}
    :=
    \frac1\ell\sum_{j=1}^{\ell}\delta_{\lambda_j(\bP_\ell)},
    \qquad
    \mu_{\mathrm{sing},\ell}^{\mathrm{sm}}
    :=
    \frac1\ell\sum_{j=1}^{\ell}\delta_{s_j(\bP_\ell)}.
\]
Then
\[
    \mu_{\mathrm{eig},\ell}^{\mathrm{sm}}\weak\delta_0,
    \qquad
    \mu_{\mathrm{sing},\ell}^{\mathrm{sm}}\weak\delta_0.
\]
More quantitatively, for every $\varepsilon>0$,
\[
    \frac1\ell\#\{j:|\lambda_j(\bP_\ell)|>\varepsilon\}
    \leq
    \frac{e^{4R^2}}{\varepsilon^2\ell},
    \qquad
    \frac1\ell\#\{j:s_j(\bP_\ell)>\varepsilon\}
    \leq
    \frac{e^{4R^2}}{\varepsilon^2\ell}.
    \label{eq:dense-softmax-lsd-collapse-count}
\]
Moreover, with the rank proxies in Definition~\ref{def:rank-proxies},
\[
    \srank(\bP_\ell)=O_R(1),
    \qquad
    \erank_1(\bP_\ell)=O_{r,R}(1),
    \qquad
    \erank_2(\bP_\ell)=O_{r,R}(1),
\]
and
\[
    \srank_\lambda(\bP_\ell)=O_{r,R}(1),
    \qquad
    \erank_\lambda(\bP_\ell)=O_{r,R}(1).
\]
\end{theorem}

\begin{proof}
The bounded feature assumption gives
$e^{-R^2}\leq e^{S_{ij}}\leq e^{R^2}$, hence
\begin{equation}
    \frac{e^{-2R^2}}{\ell}
    \leq
    P_{ij}
    \leq
    \frac{e^{2R^2}}{\ell}.
    \label{eq:dense-softmax-entry-envelope}
\end{equation}
Therefore $\|\bP_\ell\|_F^2\leq e^{4R^2}$.  Since $\bP_\ell$ is row-stochastic,
one is an eigenvalue and $\|\bP_\ell\|_{\op}\geq1$.  This proves
$\srank(\bP_\ell)=O_R(1)$.  Markov's inequality and
$\sum_j|\lambda_j(\bP_\ell)|^2\leq\sum_js_j(\bP_\ell)^2=\|\bP_\ell\|_F^2$
give \eqref{eq:dense-softmax-lsd-collapse-count}, and hence the two empirical
laws converge weakly to $\delta_0$.

It remains to control the entropy effective ranks.  Set
$A_{ij}=e^{\mathbf x_i^\top\mathbf y_j}$ and
$D_{ii}=(\sum_mA_{im})^{-1}$, so that $\bP_\ell=\bD\bA$.  For any
$0<p\leq1$,
\[
    \sum_j s_j(\bP_\ell)^p
    \leq
    \|\bD\|_{\op}^p\sum_j s_j(\bA)^p
    \leq
    \frac{e^{pR^2}}{\ell^p}\sum_j s_j(\bA)^p.
\]
Using the expansion
\[
    e^{\mathbf x_i^\top\mathbf y_j}
    =
    \sum_{\alpha\in\mathbb N^r}
    \frac{\mathbf x_i^\alpha\mathbf y_j^\alpha}{\alpha!},
\]
and the $p$-triangle inequality for Schatten quasi-norms when $p\leq1$, each
rank-one term contributes at most
$(\ell R^{2|\alpha|}/\alpha!)^p$.  Thus
\begin{equation}
    \sum_j s_j(\bP_\ell)^p
    \leq
    e^{pR^2}
    \sum_{\alpha\in\mathbb N^r}
    \frac{R^{2p|\alpha|}}{(\alpha!)^p}
    =:C_{p,r,R}<\infty.
    \label{eq:dense-softmax-schatten-p}
\end{equation}
Taking $p=1/2$ and using the R\'enyi entropy bound
$H(q)\leq H_{1/2}(q)$ for the normalized singular-value weights gives
\[
    \erank_1(\bP_\ell)
    \leq
    \frac{\left(\sum_j s_j(\bP_\ell)^{1/2}\right)^2}
    {\sum_j s_j(\bP_\ell)}
    =O_{r,R}(1),
\]
because $\sum_js_j(\bP_\ell)\geq1$.  Taking $p=1$ gives
$\|\bP_\ell\|_*=O_{r,R}(1)$, while
\eqref{eq:dense-softmax-entry-envelope} gives $\|\bP_\ell\|_F\geq e^{-2R^2}$.
The same R\'enyi bound applied to the squared singular-value weights gives
\[
    \erank_2(\bP_\ell)
    \leq
    \left(\frac{\|\bP_\ell\|_*}{\|\bP_\ell\|_F}\right)^2
    =O_{r,R}(1).
\]

For eigenvalues, Weyl's majorant theorem gives
$\sum_j|\lambda_j(\bP_\ell)|^p\leq\sum_js_j(\bP_\ell)^p$ for $p>0$.  Since
$\bP_\ell$ is row-stochastic, $\max_j|\lambda_j(\bP_\ell)|=1$.  The cases
$p=1$ and $p=1/2$ in \eqref{eq:dense-softmax-schatten-p} therefore give
$\srank_\lambda(\bP_\ell)=O_{r,R}(1)$ and, explicitly,
\[
    \erank_\lambda(\bP_\ell)
    \leq
    \frac{\left(\sum_j|\lambda_j(\bP_\ell)|^{1/2}\right)^2}
    {\sum_j|\lambda_j(\bP_\ell)|}
    =O_{r,R}(1).
\]
\end{proof}

\subsection{Causal Low-Rank Softmax}

\begin{theorem}[Causal low-rank softmax has logarithmic rank]
\label{thm:low-rank-causal-softmax-collapse}
Let
\[
    \bP_\ell
    =
    \operatorname{softmax}_{\mathrm{row}}^{\mathrm{causal}}(\bS_\ell),
    \qquad
    P_{ij}
    =
    \frac{\one_{\{j\leq i\}}e^{S_{ij}}}{\sum_{m=1}^{i}e^{S_{im}}}.
\]
Assume that the score matrix has fixed-rank bounded feature form
\begin{equation}
    S_{ij}=\mathbf x_i^\top\mathbf y_j,\qquad
    \mathbf x_i,\mathbf y_j\in\R^r,\qquad
    \max_i\|\mathbf x_i\|_2\vee\max_j\|\mathbf y_j\|_2\leq R,
    \label{eq:bounded-low-rank-causal-softmax-scores}
\end{equation}
where $r$ and $R$ do not depend on $\ell$.  Let
\[
    \mu_{\mathrm{eig},\ell}^{\mathrm{csm}}
    :=
    \frac1\ell\sum_{j=1}^{\ell}\delta_{\lambda_j(\bP_\ell)},
    \qquad
    \mu_{\mathrm{sing},\ell}^{\mathrm{csm}}
    :=
    \frac1\ell\sum_{j=1}^{\ell}\delta_{s_j(\bP_\ell)}.
\]
Then
\[
    \mu_{\mathrm{eig},\ell}^{\mathrm{csm}}\weak\delta_0,
    \qquad
    \mu_{\mathrm{sing},\ell}^{\mathrm{csm}}\weak\delta_0.
\]
More quantitatively, for every $\varepsilon>0$,
\[
    \frac1\ell\#\{j:|\lambda_j(\bP_\ell)|>\varepsilon\}
    \leq
    \frac{e^{4R^2}(1+\log\ell)}{\varepsilon^2\ell},
    \qquad
    \frac1\ell\#\{j:s_j(\bP_\ell)>\varepsilon\}
    \leq
    \frac{e^{4R^2}(1+\log\ell)}{\varepsilon^2\ell}.
    \label{eq:causal-softmax-lsd-collapse-count}
\]
Furthermore,
\[
    c_R\log\ell
    \leq
    \srank(\bP_\ell)
    \leq
    C_R(1+\log\ell),
    \qquad
    c_R\log\ell
    \leq
    \erank_1(\bP_\ell),
    \qquad
    c_R\log\ell
    \leq
    \erank_2(\bP_\ell),
    \label{eq:causal-softmax-singular-ranks-log}
\]
for constants $0<c_R<C_R<\infty$.  The diagonal eigenvalues satisfy
\[
    \frac{e^{-2R^2}}{i}
    \leq
    \lambda_i(\bP_\ell)=P_{ii}
    \leq
    \frac{e^{2R^2}}{i},
    \qquad i=1,\ldots,\ell,
\]
and hence the eigenvalue proxies obey
\[
    \srank_\lambda(\bP_\ell)
    =
    \Theta_R(\log\ell),
    \qquad
    \erank_\lambda(\bP_\ell)
    \geq
    \ell^{c_R}
    \label{eq:causal-softmax-eigenvalue-ranks}
\]
for all sufficiently large $\ell$.
\end{theorem}

\begin{proof}
The bounded feature assumption implies
\[
    e^{-R^2}\leq e^{S_{ij}}\leq e^{R^2},
    \qquad
    i e^{-R^2}
    \leq
    \sum_{m=1}^{i}e^{S_{im}}
    \leq
    i e^{R^2}.
\]
Thus
\begin{equation}
    e^{-2R^2}\frac{\one_{\{j\leq i\}}}{i}
    \leq
    P_{ij}
    \leq
    e^{2R^2}\frac{\one_{\{j\leq i\}}}{i}.
    \label{eq:causal-softmax-entry-envelope}
\end{equation}
Since $\bP_\ell$ is lower triangular, its eigenvalues are its diagonal entries,
so
\begin{equation}
    e^{-2R^2}/i
    \leq
    \lambda_i(\bP_\ell)=P_{ii}
    \leq
    e^{2R^2}/i.
    \label{eq:causal-softmax-diagonal-envelope}
\end{equation}
Moreover
\begin{equation}
    \|\bP_\ell\|_F^2
    =
    \sum_{i=1}^{\ell}\sum_{j=1}^{i}P_{ij}^2
    \asymp_R
    \sum_{i=1}^{\ell}\frac{1}{i}
    =
    \Theta_R(\log\ell).
    \label{eq:causal-softmax-frobenius-envelope}
\end{equation}

The lower bound $\|\bP_\ell\|_{\op}\geq1$ follows because $\bP_\ell$ is
row-stochastic and hence has the right eigenvector $\mathbf 1$.  For the upper
bound, let $C_{ij}=\one_{\{j\leq i\}}/i$.  By
\eqref{eq:causal-softmax-entry-envelope}, $0\leq\bP_\ell\leq e^{2R^2}\bC_\ell$
entrywise.  Since both matrices have nonnegative entries,
$|\bP_\ell x|\leq e^{2R^2}\bC_\ell|x|$ entrywise for every vector $x$, and
therefore
\[
    \|\bP_\ell\|_{\op}
    \leq
    e^{2R^2}\|\bC_\ell\|_{\op}.
\]
The Ces\`aro matrix $\bC_\ell$ has operator norm at most two by Hardy's
inequality,
\[
    \sum_{i=1}^{\ell}
    \left(\frac1i\sum_{j=1}^{i}x_j\right)^2
    \leq
    4\sum_{j=1}^{\ell}x_j^2.
\]
Therefore $\|\bP_\ell\|_{\op}=O_R(1)$.  Combining this with
\eqref{eq:causal-softmax-frobenius-envelope} gives the two-sided stable-rank
bound in \eqref{eq:causal-softmax-singular-ranks-log}.

The two counting bounds in \eqref{eq:causal-softmax-lsd-collapse-count} follow
from Markov's inequality applied to squared eigenvalue moduli and squared
singular values, using
$\sum_j|\lambda_j(\bP_\ell)|^2\leq\|\bP_\ell\|_F^2$ and
\eqref{eq:causal-softmax-frobenius-envelope}.  Since every eigenvalue and
singular value is bounded by $\|\bP_\ell\|_{\op}=O_R(1)$, the counting bounds
imply weak convergence of both empirical measures to $\delta_0$.

For $\erank_2$, apply the entropy lower bound by min-entropy:
\[
    \erank_2(\bP_\ell)
    \geq
    \frac{\|\bP_\ell\|_F^2}{\|\bP_\ell\|_{\op}^2}
    =
    \srank(\bP_\ell)
    \geq
    c_R\log\ell.
\]
The same min-entropy bound for singular-value weights gives
\[
    \erank_1(\bP_\ell)
    \geq
    \frac{\|\bP_\ell\|_*}{\|\bP_\ell\|_{\op}}
    \geq
    \frac{\|\bP_\ell\|_F^2}{\|\bP_\ell\|_{\op}^2}
    =
    \srank(\bP_\ell)
    \geq
    c_R\log\ell,
\]
where $\|\bP_\ell\|_*\geq\|\bP_\ell\|_F^2/\|\bP_\ell\|_{\op}$.

Finally, \eqref{eq:causal-softmax-diagonal-envelope} gives
\[
    \sum_i|\lambda_i(\bP_\ell)|
    =
    \Theta_R(\log\ell),
    \qquad
    \max_i|\lambda_i(\bP_\ell)|
    =
    \Theta_R(1),
\]
so $\srank_\lambda(\bP_\ell)=\Theta_R(\log\ell)$.  For the eigenvalue entropy,
the weights $p_i=|\lambda_i|/\sum_j|\lambda_j|$ are comparable, up to constants
depending only on $R$, to
\[
    \frac{1/i}{H_\ell},
    \qquad
    H_\ell:=\sum_{i=1}^{\ell}\frac1i.
\]
More precisely, there are constants $0<c_R<C_R<\infty$ such that
\[
    \frac{c_R}{iH_\ell}
    \leq p_i\leq
    \frac{C_R}{iH_\ell}.
\]
Consequently, $-\log p_i\geq \log i+\log H_\ell-\log C_R$, and therefore
\[
    -\sum_i p_i\log p_i
    \geq
    \log H_\ell
    +
    c_R
    \frac{1}{H_\ell}\sum_{i=1}^{\ell}\frac{\log i}{i}
    -
    C_R
    \geq
    c_R\log\ell
\]
for all sufficiently large $\ell$, because $H_\ell\asymp\log\ell$ and
$\sum_{i\leq\ell}(\log i)/i=(1/2+o(1))(\log\ell)^2$.  Hence
$\erank_\lambda(\bP_\ell)\geq \ell^{c_R}$.  This proves all claims.
\end{proof}

\end{document}
